\section{\texorpdfstring{Local ground space $G_L(A)$}{Local ground space G\_L(A)}}\label{sec:ground_space_parent}

Fix an MPS tensor $A=\{A^i\}_{i=0}^{d-1}$ with $A^i\in\MN{D}$ and a block
length $L\ge 1$. For a word $w=(i_1,\ldots,i_L)$, write
$A^w:=A^{i_1}\cdots A^{i_L}$. The boundary-condition parametrization is
\begin{align}
    \Gamma_L
     & \colon \MN{D}\longrightarrow(\Fin{d}^{L}\to\C),
    \label{eq:parent_gamma}                            \\
    \Gamma_L(X)(\sigma)
     & :=\tr\!\left(A^\sigma X\right).
    \notag
\end{align}
Here $A^\sigma$ abbreviates
$A^{\sigma(0)}\cdots A^{\sigma(L-1)}$. By cyclicity of the trace, this is the
same as the convention of \cite{PerezGarcia2007Matrix}, namely
$\tr(XA^\sigma)$ for the local space $\mathcal G_L^A$.

\begin{definition}[Ground-space map]\label{def:ground_space_map}
    \lean{MPSTensor.groundSpaceMap}
    \leanok
    \uses{def:eval_word}
    The map $\Gamma_L$ in (\ref{eq:parent_gamma}) is a $\C$-linear map from
    $\MN{D}$ to $(\Fin{d}^{L}\to\C)$. In tensor-network notation,
    \begin{center}
        \tenkzkernel
        % Formula/source: Papers/quant-ph_0608197/MPSarchive.tex,
        % lines 889--898, defines Gamma_L(X) with coefficients
        % tr(X A_{i_1}\cdots A_{i_L}).
        % Ink-to-index map: sigma_1,...,sigma_L are the free upper physical
        % indices; every virtual index lies on the closed A-word--X loop.
        % Contraction set: all L+1 virtual bonds around the periodic loop are summed.
        % Type verdict: a rank-L physical tensor, equivalently a length-L state.
        % Boundary signature: semantic virtual-west=0 | virtual-east=0 |
        % physical-up=L | physical-down=0; the finite schematic emits 3 up legs.
        \begin{tenkz}[rows={wire}, cols=5, boundary=periodic]
            \tn[skin=ring]{X} &
            \tn[label pos=135, ports={90:physical:$\sigma_1$}]{A} &
            \tn[label pos=135, ports={90:physical}]{A} & \tn[skin=dots]{} &
            \tn[label pos=135, ports={90:physical:$\sigma_L$}]{A}
            \tnmark[form=bracket]{(1,2) .. (1,5)}{$L$}
        \end{tenkz}
    \end{center}
    The chain denotes the word tensor $A^\sigma$ with open physical indices
    $\sigma_1,\ldots,\sigma_L$, and the red capsule denotes the boundary matrix
    $X$ inserted on the closing virtual bond before taking the trace. The brace
    records the chain length~$L$.
\end{definition}

\begin{lemma}[Evaluation formula]\label{lem:ground_space_map_apply}
    \lean{MPSTensor.groundSpaceMap_apply}
    \leanok
    \uses{def:ground_space_map}
    For every boundary matrix $X\in\MN{D}$ and every length-$L$ word $\sigma$,
    one has $\Gamma_L(X)(\sigma)=\tr(A^\sigma X)$.
\end{lemma}

\subsection{The FNW boundary convention}\label{subsec:fnw_boundary_convention}

Fannes, Nachtergaele, and Werner fix the product convention
$v(\mu_1,\ldots,\mu_N)=v(\mu_1)\cdots v(\mu_N)$ in
\cite[Equation~(5.4)]{Fannes1992Finitely} and define their boundary map in
\cite[Equation~(5.5)]{Fannes1992Finitely}.  The exact identification with the
map in Definition~\ref{def:ground_space_map} uses
$A_\mu=v(\mu)^\dagger$ and reverses the order of the physical sites.

\begin{definition}[Reversal of physical configurations]
    \label{def:physical_site_reverse_config_equiv}
    \lean{MPSTensor.physicalSiteReverseConfigEquiv}
    \leanok
    For a length-$N$ configuration $\sigma$, define
    \begin{align}
        r_N(\sigma)(j)
         & =\sigma(N-1-j),
        \label{eq:fnw_config_reverse}
    \end{align}
    where $0\leq j<N$.  This is an involutive permutation of
    $\{0,\ldots,d-1\}^{N}$.
\end{definition}

\begin{definition}[Physical-site reversal]
    \label{def:physical_site_reverse}
    \lean{MPSTensor.physicalSiteReverse}
    \leanok
    \uses{def:physical_site_reverse_config_equiv}
    Precomposition with $r_N$ defines a complex-linear isomorphism
    \begin{align}
        R_N\colon(\Fin{d}^{N}\to\C)
         & \longrightarrow(\Fin{d}^{N}\to\C).
        \label{eq:fnw_site_reverse_map}
    \end{align}
    This is a linear map of coefficient spaces and does not require an inner
    product.
\end{definition}

\begin{theorem}[Evaluation under physical-site reversal]
    \label{thm:physical_site_reverse_apply}
    \lean{MPSTensor.physicalSiteReverse_apply}
    \leanok
    \uses{def:physical_site_reverse}
    For every coefficient function $\psi$ and every configuration $\sigma$,
    \begin{align}
        (R_N\psi)(\sigma)
         & =\psi(r_N(\sigma)).
        \label{eq:fnw_site_reverse_apply}
    \end{align}
\end{theorem}

\begin{proof}
    \leanok
    \uses{def:physical_site_reverse, def:physical_site_reverse_config_equiv}
    This is the definition of precomposition by the permutation $r_N$.
\end{proof}

\begin{definition}[FNW boundary map]\label{def:fnw_boundary_map}
    \lean{MPSTensor.fnwBoundaryMap}
    \leanok
    \uses{def:eval_word}
    For $v=\{v(\mu)\}_{\mu=0}^{d-1}$, write
    $v^\sigma=v(\sigma(0))\cdots v(\sigma(N-1))$, with the empty word equal
    to $I_D$.  Define the complex-linear map
    $F_N^v\colon\MN{D}\to(\Fin{d}^{N}\to\C)$ by
    \begin{align}
        F_N^v(B)(\sigma)
         & =\tr\!\left(B(v^\sigma)^\dagger\right).
        \label{eq:fnw_boundary_map}
    \end{align}
    This is the boundary map of
    \cite[Equation~(5.5)]{Fannes1992Finitely}, with the word order fixed by
    \cite[Equation~(5.4)]{Fannes1992Finitely}.
\end{definition}

\begin{theorem}[FNW boundary-map evaluation]
    \label{thm:fnw_boundary_map_apply}
    \lean{MPSTensor.fnwBoundaryMap_apply}
    \leanok
    \uses{def:fnw_boundary_map}
    For every $B\in\MN{D}$ and every length-$N$ configuration $\sigma$,
    \begin{align}
        F_N^v(B)(\sigma)
         & =\tr\!\left(B(v^\sigma)^\dagger\right).
    \end{align}
\end{theorem}

\begin{proof}
    \leanok
    \uses{def:fnw_boundary_map}
    Evaluate the defining linear map at $B$ and then at $\sigma$.
\end{proof}

\begin{theorem}[Reversed adjoint word]
    \label{thm:fnw_boundary_map_apply_eq_trace_reverse}
    \lean{MPSTensor.fnwBoundaryMap_apply_eq_trace_reverse}
    \leanok
    \uses{thm:fnw_boundary_map_apply}
    Set $A_\mu=v(\mu)^\dagger$.  Then
    \begin{align}
        F_N^v(B)(\sigma)
         & =\tr\!\left(B A^{r_N(\sigma)}\right).
        \label{eq:fnw_boundary_reverse_word}
    \end{align}
    For $N>0$, the word on the right is
    $A_{\sigma(N-1)}\cdots A_{\sigma(0)}$.
\end{theorem}

\begin{proof}
    \leanok
    \uses{thm:fnw_boundary_map_apply}
    Conjugate transposition reverses the product:
    \begin{align}
        (v^\sigma)^\dagger
         & =A^{r_N(\sigma)}.
    \end{align}
    Thus the site order in (\ref{eq:fnw_boundary_reverse_word}) comes from the
    adjointed word itself.
\end{proof}

\begin{theorem}[FNW boundary map in the ground-space convention]
    \label{thm:fnw_boundary_map_eq_physical_site_reverse_comp_ground_space_map}
    \lean{MPSTensor.fnwBoundaryMap_eq_physicalSiteReverse_comp_groundSpaceMap}
    \leanok
    \uses{def:fnw_boundary_map, def:physical_site_reverse,
        def:ground_space_map}
    With the exact convention $A_\mu=v(\mu)^\dagger$, one has
    \begin{align}
        F_N^v
         & =R_N\circ\Gamma_N^A.
        \label{eq:fnw_boundary_ground_space_map}
    \end{align}
\end{theorem}

\begin{proof}
    \leanok
    \uses{thm:fnw_boundary_map_apply_eq_trace_reverse,
        thm:physical_site_reverse_apply, lem:ground_space_map_apply}
    For every $B\in\MN{D}$ and every configuration $\sigma$,
    \begin{align}
        (R_N\Gamma_N^A(B))(\sigma)
         & =\Gamma_N^A(B)(r_N(\sigma))          \\
         & =\tr\!\left(A^{r_N(\sigma)}B\right)  \\
         & =\tr\!\left(B A^{r_N(\sigma)}\right) \\
         & =F_N^v(B)(\sigma).
    \end{align}
    The ordered list associated with $r_N(\sigma)$ is the reverse of the list
    associated with $\sigma$, as in (\ref{eq:fnw_config_reverse}).  Together
    with the adjointed word in
    Theorem~\ref{thm:fnw_boundary_map_apply_eq_trace_reverse}, this gives the
    physical-site reversal.  Trace cyclicity is used only in the third
    equality, where it moves the boundary matrix $B$ across the word.
\end{proof}

The maps in this subsection are raw algebraic maps on matrices and coefficient
functions.  The $\rho$-weighted Hilbert-space structure of
\cite[Equation~(5.6)]{Fannes1992Finitely} is a separate construction; no
choice of $\rho$ enters (\ref{eq:fnw_boundary_ground_space_map}).

\begin{definition}[Local ground space]\label{def:ground_space_parent}
    \lean{MPSTensor.groundSpace}
    \leanok
    \uses{def:ground_space_map}
    The \emph{local ground space} is the image
    $G_L(A):=\range(\Gamma_L)\subseteq(\Fin{d}^{L}\to\C)$.
\end{definition}

\begin{lemma}[Dimension bound]\label{lem:ground_space_finrank_le}
    \lean{MPSTensor.groundSpace_finrank_le}
    \leanok
    \uses{def:ground_space_parent}
    For every $L$, one has $\dim G_L(A)\le D^2$.
\end{lemma}

\begin{proof}
    \leanok
    Since $G_L(A)$ is the range of a linear map from $\MN{D}$, one has
    $\dim G_L(A)\le\dim\MN{D}=D^2$.
\end{proof}

\begin{lemma}[Ambient-space dimension]\label{lem:nsite_space_finrank}
    \lean{MPSTensor.nSiteSpace_finrank}
    \leanok
    The above coefficient space satisfies
    $\dim(\Fin{d}^{L}\to\C)=d^L$.
\end{lemma}

\begin{lemma}[Nontriviality criterion]\label{lem:ground_space_ne_top}
    \lean{MPSTensor.groundSpace_ne_top}
    \leanok
    \uses{lem:ground_space_finrank_le, lem:nsite_space_finrank}
    If $d^L>D^2$, then $G_L(A)$ is a proper subspace of the local Hilbert
    space, namely $G_L(A)\neq(\Fin{d}^{L}\to\C)$.
\end{lemma}

\begin{proof}
    \leanok
    If $G_L(A)$ were the whole space, its dimension would be $d^L$ by
    Lemma~\ref{lem:nsite_space_finrank}; Lemma~\ref{lem:ground_space_finrank_le}
    would then give the contradiction $d^L\le D^2$.
\end{proof}

\section{Parent interaction and chain Hamiltonian}\label{sec:parent_interaction_defs}

The canonical parent interaction at range $L$ is
\begin{align}
    h_L(A)
     & =\Id-\Pi_{G_L(A)}
    =\Pi_{G_L(A)^\perp},
    \label{eq:parent_interaction} \\
    \ker h_L(A)
     & =G_L(A).
    \notag
\end{align}
The sources also allow any positive local interaction with this kernel; the
orthogonal projector is the canonical representative used here.

\begin{definition}[Local ground space under the canonical \texorpdfstring{$\ell^2$}{l2} identification]\label{def:ground_space_es}
    \lean{MPSTensor.groundSpaceES}
    \leanok
    \uses{def:nsite_space_parent, def:ground_space_parent}
    Let $\iota_L\colon(\Fin{d}^{L}\to\C)\simeq\ell^2(\Fin{d}^{L})$ be the
    canonical computational-basis identification. The local ground space under
    this $\ell^2$ identification is
    $G_L^{\mathrm{ES}}(A):=\iota_L(G_L(A))
        \subseteq\ell^2(\Fin{d}^{L})$.
\end{definition}

\begin{lemma}[Membership under the canonical \texorpdfstring{$\ell^2$}{l2} identification]\label{lem:mem_ground_space_es}
    \lean{MPSTensor.mem_groundSpaceES_iff}
    \leanok
    \uses{def:ground_space_es, def:ground_space_parent}
    For $v\in\ell^2(\Fin{d}^{L})$, one has $v\in G_L^{\mathrm{ES}}(A)$ if and
    only if $\iota_L^{-1}v\in G_L(A)$.
\end{lemma}

\begin{proof}
    \leanok
    This is the defining equation
    $G_L^{\mathrm{ES}}(A)=\iota_L(G_L(A))$.
\end{proof}

\begin{definition}[Hilbert-space boundary map]
    \label{def:ground_space_map_es}
    \lean{MPSTensor.groundSpaceMapES}
    \leanok
    \uses{def:ground_space_map}
    Let $U\colon\MN{D}\to\ell^2(\Fin{D}\times\Fin{D})$ be Frobenius
    vectorization in column--row coordinates.  The Hilbert-space boundary map
    is
    \begin{align}
        \Gamma_L^{\mathrm{ES}}
         & :=\iota_L\circ\Gamma_L\circ U^{-1}
        \colon\ell^2(\Fin{D}\times\Fin{D})
        \longrightarrow\ell^2(\Fin{d}^{L}).
    \end{align}
    Thus $\Gamma_L^{\mathrm{ES}}$ has a Euclidean coordinate-space domain,
    whereas the original map $\Gamma_L$ has matrix domain $\MN{D}$.
\end{definition}

\begin{theorem}[Boundary-map transport under Frobenius vectorization]
    \label{thm:ground_space_map_es_frobenius}
    \lean{MPSTensor.groundSpaceMapES_frobeniusEquivEuclidean_apply}
    \leanok
    \uses{def:ground_space_map_es}
    Let $U\colon\MN{D}\to\ell^2(\Fin{D}\times\Fin{D})$ be Frobenius
    vectorization, with the target indexed in column--row order, and let
    $\iota_L$ be the computational-basis identification on the physical
    coefficient space.  Then, for every $X\in\MN{D}$,
    $\Gamma_L^{\mathrm{ES}}(U(X))=\iota_L(\Gamma_L(X))$.
\end{theorem}

\begin{proof}
    \leanok
    \uses{def:ground_space_map}
    By definition, $\Gamma_L^{\mathrm{ES}}$ is the composite of $U^{-1}$,
    the original boundary map $\Gamma_L$, and $\iota_L$.  Applying this
    composite to $U(X)$ and cancelling $U^{-1}U$ gives the claimed identity.
\end{proof}

\begin{theorem}[Boundary map on a matrix unit]
    \label{thm:ground_space_map_es_single}
    \lean{MPSTensor.groundSpaceMapES_single}
    \leanok
    \uses{def:ground_space_map_es}
    Let $e_{(b,a)}$ denote the Frobenius coordinate vector representing the
    virtual matrix unit $E_{ab}$.  Then
    $\Gamma_L^{\mathrm{ES}}(e_{(b,a)})=\iota_L(\Gamma_L(E_{ab}))$.
\end{theorem}

\begin{proof}
    \leanok
    \uses{thm:ground_space_map_es_frobenius}
    Frobenius vectorization records a matrix entry $(a,b)$ at the
    column--row coordinate $(b,a)$, so $U(E_{ab})=e_{(b,a)}$.  Substitute
    this identity into the boundary-map transport theorem.
\end{proof}

\begin{theorem}[Exact range of the Hilbert-space boundary map]
    \label{thm:range_ground_space_map_es}
    \lean{MPSTensor.range_groundSpaceMapES}
    \leanok
    \uses{def:ground_space_map_es, def:ground_space_es}
    The range of $\Gamma_L^{\mathrm{ES}}$ is exactly
    $G_L^{\mathrm{ES}}(A)$.
\end{theorem}

\begin{proof}
    \leanok
    \uses{thm:ground_space_map_es_frobenius, def:ground_space_parent}
    Every virtual vector is $U(X)$ for a unique matrix $X$, and the transport
    identity sends it to $\iota_L(\Gamma_L(X))$.  These are exactly the
    vectors in $\iota_L(\range\Gamma_L)=G_L^{\mathrm{ES}}(A)$, proving both
    inclusions.
\end{proof}

\begin{theorem}[Hilbert-space boundary-map injectivity]
    \label{thm:ground_space_map_es_injective_block}
    \lean{MPSTensor.groundSpaceMapES_injective_of_isNBlkInjective}
    \leanok
    \uses{def:ground_space_map_es, def:l_blk_injective}
    If $A$ is $L$-block injective, then
    $\Gamma_L^{\mathrm{ES}}$ is injective.
\end{theorem}

\begin{proof}
    \leanok
    \uses{thm:ground_space_map_es_frobenius, thm:ground_space_map_injective_blk}
    Write each virtual vector uniquely as $U(X)$ by Frobenius vectorization.
    The transport identity reduces equality of their images under
    $\Gamma_L^{\mathrm{ES}}$ to equality of the images of the corresponding
    matrices under $\Gamma_L$.  Block injectivity makes the latter map
    injective.
\end{proof}

\begin{theorem}[Long-length Hilbert-space boundary-map injectivity]
    \label{thm:ground_space_map_es_injective_block_of_le}
    \lean{MPSTensor.groundSpaceMapES_injective_of_isNBlkInjective_of_le}
    \leanok
    \uses{def:ground_space_map_es, def:l_blk_injective}
    Suppose that $L_0>0$, that $A$ is $L_0$-block injective, and that
    $L\ge L_0$.  Then $\Gamma_L^{\mathrm{ES}}$ is injective.
\end{theorem}

\begin{proof}
    \leanok
    \uses{thm:ground_space_map_es_injective_block, thm:wordspan_blk_inj_succ}
    Positive-length propagation makes $A$ block injective at length $L$.
    The preceding theorem then gives injectivity of
    $\Gamma_L^{\mathrm{ES}}$.
\end{proof}

\begin{definition}[Ground-space Gram operator]
    \label{def:ground_space_gram}
    \lean{MPSTensor.groundSpaceGram}
    \leanok
    \uses{def:ground_space_map_es}
    The boundary Gram operator is
    \begin{align}
        \mathcal K_L
        :=(\Gamma_L^{\mathrm{ES}})^\dagger\Gamma_L^{\mathrm{ES}}.
    \end{align}
\end{definition}

The FNW convention starts from an isometry
$V\colon\C^D\to\C^d\otimes\C^D$ written as
$V\xi=\sum_\mu\ket{\mu}\otimes v(\mu)^\dagger\xi$ in
\cite[Equation~(5.2)]{Fannes1992Finitely}. Equations~(5.3b) and~(5.3d)
of that paper give
$\sum_\mu v(\mu)v(\mu)^\dagger=I_D$ and
$E(B)=\sum_\mu v(\mu)Bv(\mu)^\dagger$. Under the identification
$A^\mu=v(\mu)^\dagger$, this map uses the adjoint tensor letters. The
asymptotic identity
$E^n(B)\to\tr(\rho B)I_D$ in
\cite[Equation~(5.1)]{Fannes1992Finitely} is a separate assertion and is not
used in the convention identities below.

\begin{definition}[FNW transfer map]
    \label{def:fnw_transfer_map}
    \lean{MPSTensor.fnwTransferMap}
    \leanok
    \uses{def:transfer_map}
    Let $A$ be an MPS tensor, with $A^\mu=v(\mu)^\dagger$ in the notation of
    \cite[Equations~(5.2)--(5.3)]{Fannes1992Finitely}. Its FNW transfer map is
    \begin{align}
        E_{\mathrm{FNW},A}(X)
         & :=\sum_{\mu=0}^{d-1}(A^\mu)^\dagger X A^\mu
        =\sum_{\mu=0}^{d-1}v(\mu)Xv(\mu)^\dagger.
        \label{eq:fnw_transfer_map}
    \end{align}
\end{definition}

\begin{theorem}[FNW map as the trace-pairing adjoint]
    \label{thm:fnw_transfer_map_trace_adjoint}
    \lean{MPSTensor.fnwTransferMap_eq_traceAdjointMap}
    \leanok
    \uses{def:fnw_transfer_map, def:transfer_map}
    Write $\E_A^\sharp$ for the adjoint of $\E_A$ with respect to the
    bilinear trace pairing, characterized by
    \begin{align}
        \tr\!\left(\E_A(X)Y\right)
         & =\tr\!\left(X\E_A^\sharp(Y)\right).
        \label{eq:fnw_transfer_map_trace_adjoint}
    \end{align}
    Then $E_{\mathrm{FNW},A}=\E_A^\sharp$. Here $\sharp$ denotes the adjoint
    for $(X,Y)\mapsto\tr(XY)$, so the theorem is phrased as a bilinear
    trace-adjoint identity. For the unweighted Frobenius inner product, the
    same map also agrees with the Hilbert-space adjoint of $\E_A$. This
    Frobenius adjoint is generally different from the adjoint for the
    $\rho$-weighted inner product in
    \cite[Equation~(5.6)]{Fannes1992Finitely}.
\end{theorem}

\begin{proof}
    \leanok
    \uses{def:fnw_transfer_map, def:transfer_map}
    Expanding both maps and using cyclicity of trace gives
    \begin{align}
        \tr\!\left(\E_A(X)Y\right)
         & =\sum_\mu\tr\!\left(A^\mu X(A^\mu)^\dagger Y\right)
        =\sum_\mu\tr\!\left(X(A^\mu)^\dagger YA^\mu\right).
    \end{align}
    The final sum is $\tr(XE_{\mathrm{FNW},A}(Y))$.
\end{proof}

\begin{theorem}[Unitality of the FNW transfer map]
    \label{thm:fnw_transfer_map_one}
    \lean{MPSTensor.fnwTransferMap_one}
    \leanok
    \uses{def:fnw_transfer_map, def:left_canonical_tensor}
    If $A$ is left canonical, then
    \begin{align}
        E_{\mathrm{FNW},A}(I_D) & =I_D.
        \label{eq:fnw_transfer_map_one}
    \end{align}
    Under $A^\mu=v(\mu)^\dagger$, this is the normalization
    $\sum_\mu v(\mu)v(\mu)^\dagger=I_D$ from
    \cite[Equation~(5.3b)]{Fannes1992Finitely}.
\end{theorem}

\begin{proof}
    \leanok
    \uses{thm:fnw_transfer_map_trace_adjoint, lem:trace_transfer_map_left_canonical}
    Left canonicity makes $\E_A$ trace-preserving. Therefore, for every $X$,
    \begin{align}
        \tr\!\left(XE_{\mathrm{FNW},A}(I_D)\right)
         & =\tr\!\left(\E_A(X)\right)=\tr(X).
    \end{align}
    Nondegeneracy of the trace pairing gives
    $E_{\mathrm{FNW},A}(I_D)=I_D$.
\end{proof}

\begin{theorem}[FNW transfer pairing identity]
    \label{thm:trace_mul_fnw_transfer_map}
    \lean{MPSTensor.trace_mul_fnwTransferMap}
    \leanok
    \uses{def:fnw_transfer_map, def:transfer_map}
    For all $\rho,X\in\MN{D}$,
    \begin{align}
        \tr\!\left(\rho E_{\mathrm{FNW},A}(X)\right)
         & =\tr\!\left(\E_A(\rho)X\right).
        \label{eq:trace_mul_fnw_transfer_map}
    \end{align}
    This is the explicit trace-pairing form of the convention
    $A^\mu=v(\mu)^\dagger$ for the map in
    \cite[Equations~(5.2)--(5.3)]{Fannes1992Finitely}.
\end{theorem}

\begin{proof}
    \leanok
    \uses{thm:fnw_transfer_map_trace_adjoint}
    Substitute the trace-pairing adjoint identity and use cyclicity of trace.
\end{proof}

These identities only compare the two transfer-map conventions. They do not
include the weighted decay estimate of
\cite[Lemma~5.2]{Fannes1992Finitely}, and they make no assertion of a
constant $C_\lambda=k^2$.

\begin{definition}[Normalized limiting FNW observable map]
    \label{def:fnw_limit_map}
    \lean{MPSTensor.fnwLimitMap}
    \lean{MPSTensor.fnwLimitMap_apply}
    \lean{MPSTensor.fnwLimitMap_apply_of_trace_eq_one}
    \leanok
    Let $\rho\in\MN{D}$ satisfy $\tr(\rho)\neq0$. The normalized rank-one map
    associated with $\rho$ is
    \begin{align}
        E_\infty^\rho(B)
         & :=\frac{\tr(\!\rho B)}{\tr(\rho)}I_D.
        \label{eq:fnw_limit_map}
    \end{align}
    If $\tr(\rho)=1$, this becomes
    $E_\infty^\rho(B)=\tr(\rho B)I_D$, which is the limiting map in
    \cite[Equation~(5.1)]{Fannes1992Finitely}.
\end{definition}

\begin{theorem}[Limiting map as a trace-pairing adjoint]
    \label{thm:fnw_limit_map_trace_adjoint}
    \lean{MPSTensor.fnwLimitMap_eq_traceAdjointMap_fixedPointProj}
    \leanok
    \uses{def:fnw_limit_map}
    Define the fixed-point projection
    \begin{align}
        P_\rho(X)
         & :=\frac{\tr(X)}{\tr(\rho)}\rho.
        \label{eq:fnw_fixed_point_projection}
    \end{align}
    Then $E_\infty^\rho=P_\rho^\sharp$, where $\sharp$ denotes the adjoint for
    the bilinear trace pairing. For the trace-one Hermitian density matrix used
    by FNW, the same map also agrees with the Hilbert-space adjoint for the
    unweighted Frobenius inner product. This Frobenius adjoint is generally
    different from the adjoint for the $\rho$-weighted inner product of
    \cite[Equation~(5.6)]{Fannes1992Finitely}.
\end{theorem}

\begin{proof}
    \leanok
    \uses{def:fnw_limit_map}
    For all $X,B\in\MN{D}$,
    \begin{align}
        \tr(P_\rho(X)B)
         & =\frac{\tr(X)\tr(\rho B)}{\tr(\rho)}
        =\tr\!\left(XE_\infty^\rho(B)\right).
    \end{align}
    Nondegeneracy of the trace pairing proves the equality of maps.
\end{proof}

\begin{theorem}[Pointwise idempotence of the limiting map]
    \label{thm:fnw_limit_map_idempotent}
    \lean{MPSTensor.fnwLimitMap_idempotent}
    \leanok
    \uses{def:fnw_limit_map}
    For every $B\in\MN{D}$,
    \begin{align}
        E_\infty^\rho\!\left(E_\infty^\rho(B)\right)
         & =E_\infty^\rho(B).
        \label{eq:fnw_limit_map_idempotent}
    \end{align}
\end{theorem}

\begin{proof}
    \leanok
    \uses{def:fnw_limit_map}
    Substitute~(\ref{eq:fnw_limit_map}) and use
    $\tr(\rho I_D)=\tr(\rho)$. The nonzero trace cancels the normalizing
    denominator.
\end{proof}

\begin{theorem}[The limiting map is a projection]
    \label{thm:fnw_limit_map_projection}
    \lean{MPSTensor.fnwLimitMap_mul_self}
    \leanok
    \uses{def:fnw_limit_map}
    As an endomorphism of $\MN{D}$, the limiting map satisfies
    \begin{align}
        (E_\infty^\rho)^2 & =E_\infty^\rho.
        \label{eq:fnw_limit_map_projection}
    \end{align}
\end{theorem}

\begin{proof}
    \leanok
    \uses{thm:fnw_limit_map_idempotent}
    Apply Theorem~\ref{thm:fnw_limit_map_idempotent} to every matrix $B$.
\end{proof}

\begin{theorem}[Unital absorption of the limiting map]
    \label{thm:fnw_transfer_absorbs_limit}
    \lean{MPSTensor.fnwTransferMap_mul_fnwLimitMap}
    \leanok
    \uses{def:fnw_limit_map, def:transfer_map, def:left_canonical_tensor}
    If $A$ is left canonical, then the FNW observable map in
    (\ref{eq:fnw_transfer_map}) satisfies
    \begin{align}
        E_{\mathrm{FNW},A}E_\infty^\rho & =E_\infty^\rho.
        \label{eq:fnw_transfer_absorbs_limit}
    \end{align}
\end{theorem}

\begin{proof}
    \leanok
    \uses{def:fnw_limit_map, def:left_canonical_tensor, thm:fnw_transfer_map_one}
    Apply Theorem~\ref{thm:fnw_transfer_map_one} to obtain
    $E_{\mathrm{FNW},A}(I_D)=\sum_\mu(A^\mu)^\dagger A^\mu=I_D$.
    Hence
    \begin{align}
        E_{\mathrm{FNW},A}\!\left(E_\infty^\rho(B)\right)
         & =\frac{\tr(\rho B)}{\tr(\rho)}E_{\mathrm{FNW},A}(I_D)
        =E_\infty^\rho(B).
    \end{align}
\end{proof}

\begin{theorem}[Stationary absorption by the limiting map]
    \label{thm:fnw_limit_absorbs_transfer}
    \lean{MPSTensor.fnwLimitMap_mul_fnwTransferMap}
    \leanok
    \uses{def:fnw_limit_map, def:transfer_map}
    Suppose that $\rho$ is stationary for the Schrödinger transfer map,
    $\E_A(\rho)=\rho$. Then
    \begin{align}
        E_\infty^\rho E_{\mathrm{FNW},A} & =E_\infty^\rho.
        \label{eq:fnw_limit_absorbs_transfer}
    \end{align}
\end{theorem}

\begin{proof}
    \leanok
    \uses{def:fnw_limit_map, thm:trace_mul_fnw_transfer_map}
    Apply Theorem~\ref{thm:trace_mul_fnw_transfer_map} and stationarity to get
    \begin{align}
        \tr\!\left(\rho E_{\mathrm{FNW},A}(B)\right)
         & =\tr\!\left(\E_A(\rho)B\right)=\tr(\rho B).
    \end{align}
    Substitute this equality into~(\ref{eq:fnw_limit_map}).
\end{proof}

\begin{theorem}[Positive powers of the FNW remainder]
    \label{thm:fnw_remainder_power}
    \lean{MPSTensor.fnwTransferMap_sub_fnwLimitMap_pow}
    \leanok
    \uses{def:fnw_limit_map, def:transfer_map, def:left_canonical_tensor}
    Suppose that $A$ is left canonical and that
    $\E_A(\rho)=\rho$. For every integer $n\geq1$,
    \begin{align}
        (E_{\mathrm{FNW},A}-E_\infty^\rho)^n
         & =E_{\mathrm{FNW},A}^n-E_\infty^\rho.
        \label{eq:fnw_remainder_power}
    \end{align}
    The restriction $n\geq1$ is essential. At $n=0$, the two sides would be
    $I$ and $I-E_\infty^\rho$, respectively, and
    $E_\infty^\rho(I_D)=I_D$.
\end{theorem}

\begin{proof}
    \leanok
    \uses{thm:fnw_limit_map_projection, thm:fnw_transfer_absorbs_limit, thm:fnw_limit_absorbs_transfer}
    Write $E=E_{\mathrm{FNW},A}$ and $P=E_\infty^\rho$. The three preceding
    identities give
    $EP=PE=P^2=P$, and hence $E^mP=P$ for every $m\geq0$. The case $n=1$ is
    immediate. If $(E-P)^n=E^n-P$, then
    \begin{align}
        (E-P)^{n+1}
         & =(E^n-P)(E-P)
        =E^{n+1}-E^nP-PE+P^2
        =E^{n+1}-P.
    \end{align}
\end{proof}

\begin{theorem}[Primitive-MPS remainder identity]
    \label{thm:primitive_fnw_remainder_power}
    \lean{MPSTensor.IsPrimitiveMPS.fnwTransferMap_sub_fnwLimitMap_pow}
    \leanok
    \uses{def:primitive_mps, def:fnw_limit_map}
    Let $A$ be a primitive MPS tensor of nonzero bond dimension with fixed-point
    witness $\rho$. Then, for every integer $n\geq1$,
    \begin{align}
        (E_{\mathrm{FNW},A}-E_\infty^\rho)^n
         & =E_{\mathrm{FNW},A}^n-E_\infty^\rho.
    \end{align}
\end{theorem}

\begin{proof}
    \leanok
    \uses{thm:fnw_remainder_power}
    Primitivity supplies $\tr(\rho)\neq0$, left canonicity, and the stationary
    equation $\E_A(\rho)=\rho$. Apply
    Theorem~\ref{thm:fnw_remainder_power}.
\end{proof}

These projection and absorption identities precede the weighted Hilbert-space
estimate in \cite[Lemma~5.2]{Fannes1992Finitely}. They do not assert a
$\rho$-weighted norm bound or any quantitative decay constant, and they do not
include the later conclusions of Lemmas~5.3 and~6.2 of that paper.

\begin{definition}[Operator-level Gram reshuffling]
    \label{def:gram_reshuffle}
    \lean{Matrix.gramReshuffle}
    \lean{Matrix.gramReshuffleMatrix}
    \leanok

    For a complex-linear map $\mathcal L\colon\MN{D}\to\MN{D}$, let
    $R(\mathcal L)$ be the matrix with entries
    \begin{align}
        R(\mathcal L)_{(b,a),(e,c)}
         & :=J(\mathcal L)_{(c,e),(a,b)}.
        \label{eq:gram_reshuffle_matrix_entries}
    \end{align}
    The operator $\mathscr R(\mathcal L)$ on
    $\ell^2(\Fin{D}\times\Fin{D})$ is represented by the matrix $R(\mathcal L)$
    in the standard orthonormal basis.
    Thus the row--column pair $((b,a),(e,c))$ is sent explicitly to the Choi
    row--column pair $((c,e),(a,b))$.
\end{definition}

\begin{theorem}[Gram reshuffling commutes with subtraction]
    \label{thm:gram_reshuffle_sub}
    \lean{Matrix.gramReshuffle_sub}
    \leanok
    \uses{def:gram_reshuffle}
    For complex-linear maps
    $\mathcal L,\mathcal M\colon\MN{D}\to\MN{D}$,
    \begin{align}
        \mathscr R(\mathcal L-\mathcal M)
         & =\mathscr R(\mathcal L)-\mathscr R(\mathcal M).
        \label{eq:gram_reshuffle_sub}
    \end{align}
\end{theorem}

\begin{proof}
    \leanok
    \uses{def:gram_reshuffle}
    The defining matrix entries of the Choi reshuffling are linear in the map.
\end{proof}

\begin{theorem}[Matrix-unit coordinates of Gram reshuffling]
    \label{thm:gram_reshuffle_single}
    \lean{Matrix.inner_single_gramReshuffle_single}
    \leanok
    \uses{def:gram_reshuffle}
    Let $\mathcal L\colon\MN{D}\to\MN{D}$ be complex-linear and let
    $\mathbf e_{(j,i)}$ be the standard unit vector at $(j,i)$.  For all
    $a,b,c,e\in\Fin{D}$,
    \begin{align}
        \bigl\langle \mathbf e_{(b,a)},
        \mathscr R(\mathcal L)\mathbf e_{(e,c)}\bigr\rangle_{\ell^2}
         & =J(\mathcal L)_{(c,e),(a,b)}.
        \label{eq:gram_reshuffle_single}
    \end{align}
    This is the coordinate identity in
    (\ref{eq:gram_reshuffle_matrix_entries}); it is not a Hilbert--Schmidt pairing of
    $E_{ab}$ with $\mathcal L(E_{ce})$.
\end{theorem}

\begin{proof}
    \leanok
    \uses{def:gram_reshuffle}
    Evaluate the matrix of $\mathscr R(\mathcal L)$ on the two standard basis
    vectors.
\end{proof}

\begin{theorem}[Euclidean operator norm bounded by Frobenius mass]
    \label{thm:euclidean_operator_norm_sq_le_frobenius}
    \lean{Matrix.toEuclideanCLM_norm_sq_le_sum_norm_sq}
    \leanok
    Let $I$ be a finite index set and let $B=(B_{ij})_{i,j\in I}$ be a complex
    matrix.  Then the squared operator norm of $B$ on $\ell^2(I)$ satisfies
    \begin{align}
        \lVert B\rVert_{\ell^2(I)\to\ell^2(I)}^2
         & \leq \sum_{(i,j)\in I\times I}\lvert B_{ij}\rvert^2.
        \label{eq:euclidean_operator_norm_sq_le_frobenius}
    \end{align}
\end{theorem}

\begin{proof}
    \leanok
    Apply Cauchy--Schwarz to each row of $Bx$, sum over the rows, and take the
    supremum over unit vectors $x\in\ell^2(I)$.
\end{proof}

\begin{theorem}[Squared mass of one column]
    \label{thm:sum_column_norm_sq_le_l2_operator_norm_sq}
    \lean{Matrix.sum_column_norm_sq_le_norm_sq}
    \leanok
    Let $I$ be a finite index set, let $A=(A_{ij})_{i,j\in I}$ be a complex
    matrix, and let $j\in I$.  Then
    \begin{align}
        \sum_{i\in I}\lvert A_{ij}\rvert^2
         & \leq \lVert A\rVert_{\ell^2(I)\to\ell^2(I)}^2.
        \label{eq:sum_column_norm_sq_le_l2_operator_norm_sq}
    \end{align}
\end{theorem}

\begin{proof}
    \leanok
    Apply $A$ to the unit vector supported at $j$.  Its image is the $j$th
    column of $A$, while the vector itself has Euclidean norm one.
\end{proof}

\begin{theorem}[Entrywise mass bounded by the $L^2$ operator norm]
    \label{thm:sum_entry_norm_sq_le_card_mul_l2_operator_norm_sq}
    \lean{Matrix.sum_entry_norm_sq_le_card_mul_norm_sq}
    \leanok
    Let $I$ be a finite index set and let $A=(A_{ij})_{i,j\in I}$ be a complex
    matrix.  Then
    \begin{align}
        \sum_{(i,j)\in I\times I}\lvert A_{ij}\rvert^2
         & \leq |I|\,
        \lVert A\rVert_{\ell^2(I)\to\ell^2(I)}^2.
        \label{eq:sum_entry_norm_sq_le_card_mul_l2_operator_norm_sq}
    \end{align}
\end{theorem}

\begin{proof}
    \leanok
    \uses{thm:sum_column_norm_sq_le_l2_operator_norm_sq}
    Sum Theorem~\ref{thm:sum_column_norm_sq_le_l2_operator_norm_sq} over all
    columns.
\end{proof}

\begin{theorem}[$L^2$ operator norm of a matrix unit]
    \label{thm:l2_operator_norm_matrix_unit}
    \lean{Matrix.l2_opNorm_single_one}
    \leanok
    Let $I$ be a finite index set.  For every $i,j\in I$, the matrix unit
    $E_{ij}$ satisfies
    \begin{align}
        \lVert E_{ij}\rVert_{\ell^2(I)\to\ell^2(I)}=1.
        \label{eq:l2_operator_norm_matrix_unit}
    \end{align}
\end{theorem}

\begin{proof}
    \leanok
    The matrix unit sends the unit vector at $j$ to the unit vector at $i$ and
    sends every vector to a scalar multiple of that unit vector.
\end{proof}

\begin{theorem}[Entrywise Frobenius invariance of Gram reshuffling]
    \label{thm:gram_reshuffle_matrix_norm_sq}
    \lean{Matrix.gramReshuffleMatrix_norm_sq_eq_rectangularChoi_norm_sq}
    \leanok
    \uses{def:gram_reshuffle}
    For every complex-linear map $\mathcal L\colon\MN{D}\to\MN{D}$, the sum
    of squared entry norms of $R(\mathcal L)$ equals that of the rectangular
    Choi matrix:
    \begin{align}
        \sum_{p:(\Fin{D}\times\Fin{D})\times(\Fin{D}\times\Fin{D})}
        \bigl\|R(\mathcal L)_{p_1,p_2}\bigr\|^2
         & =
        \sum_{p:(\Fin{D}\times\Fin{D})\times(\Fin{D}\times\Fin{D})}
        \bigl\|J(\mathcal L)_{p_1,p_2}\bigr\|^2.
        \label{eq:gram_reshuffle_matrix_norm_sq}
    \end{align}
    The two matrices differ only by a reshuffling of indices; the equality is
    entrywise Frobenius, not operator norm.
\end{theorem}

\begin{proof}
    \leanok
    \uses{def:gram_reshuffle}
    The map $((b,a),(e,c))\mapsto((c,e),(a,b))$ is an involution (bijection)
    on the product index space.  Reindexing the sum of squared norms by this
    involution yields the claimed equality.
\end{proof}

\begin{theorem}[Operator bound for a reshuffled Choi matrix]
    \label{thm:gram_reshuffle_norm_sq_le_rectangular_choi}
    \lean{Matrix.gramReshuffle_norm_sq_le_rectangularChoi_norm_sq}
    \leanok
    \uses{def:gram_reshuffle}
    Let $I$ be a finite index set and let
    $\mathcal L\colon\mathbb C^{I\times I}\to\mathbb C^{I\times I}$ be
    complex-linear.  Then
    \begin{align}
        \bigl\lVert\mathscr R(\mathcal L)\bigr\rVert_{
            \ell^2(I\times I)\to\ell^2(I\times I)}^2
         & \leq
        \sum_{u,v\in I\times I}
        \bigl\lvert J(\mathcal L)_{u,v}\bigr\rvert^2.
        \label{eq:gram_reshuffle_norm_sq_le_rectangular_choi}
    \end{align}
\end{theorem}

\begin{proof}
    \leanok
    \uses{def:gram_reshuffle, thm:euclidean_operator_norm_sq_le_frobenius, thm:gram_reshuffle_matrix_norm_sq}
    Apply Theorem~\ref{thm:euclidean_operator_norm_sq_le_frobenius} to
    $R(\mathcal L)$, then use
    Theorem~\ref{thm:gram_reshuffle_matrix_norm_sq} to rewrite its entrywise
    Frobenius mass as that of $J(\mathcal L)$.
\end{proof}

\begin{theorem}[Dimension bound for Gram reshuffling]
    \label{thm:gram_reshuffle_norm_sq_le_card_cube_mul_operator_norm_sq}
    \lean{Matrix.gramReshuffle_norm_sq_le_card_cube_mul_opNorm_sq}
    \leanok
    \uses{def:gram_reshuffle}
    Let $I$ be a finite index set, equip $\mathbb C^{I\times I}$ with its
    $L^2$ operator norm, and let
    $\mathcal L\colon\mathbb C^{I\times I}\to\mathbb C^{I\times I}$ be
    complex-linear.  Denote the operator norm of $\mathcal L$ for these matrix
    norms by $\lVert\mathcal L\rVert_{\mathrm{op},L^2}$.  Then
    \begin{align}
        \bigl\lVert\mathscr R(\mathcal L)\bigr\rVert_{
            \ell^2(I\times I)\to\ell^2(I\times I)}^2
         & \leq |I|^3\,
        \lVert\mathcal L\rVert_{\mathrm{op},L^2}^2.
        \label{eq:gram_reshuffle_norm_sq_le_card_cube_mul_operator_norm_sq}
    \end{align}
\end{theorem}

\begin{proof}
    \leanok
    \uses{thm:gram_reshuffle_norm_sq_le_rectangular_choi, thm:sum_entry_norm_sq_le_card_mul_l2_operator_norm_sq, thm:l2_operator_norm_matrix_unit}
    Use Theorem~\ref{thm:gram_reshuffle_norm_sq_le_rectangular_choi} and expand
    the Choi matrix in matrix units.  For each of the $|I|^2$ input matrix
    units, Theorem~\ref{thm:sum_entry_norm_sq_le_card_mul_l2_operator_norm_sq}
    contributes a factor $|I|$.  The operator-norm estimate for $\mathcal L$
    and Theorem~\ref{thm:l2_operator_norm_matrix_unit} bound the remaining
    matrix norm.
\end{proof}

\begin{theorem}[Geometric power bound at a prescribed rate]
    \label{thm:geometric_bound_of_spectral_radius_lt}
    \lean{geometric_bound_of_spectralRadius_lt}
    \leanok
    Let $A$ be a complex Banach algebra, let $a\in A$, and let
    $\lambda\geq0$.  If $\rho_{\mathrm{spec}}(a)<\lambda$, then there is a
    constant $C_\lambda>0$ such that
    \begin{align}
        \lVert a^n\rVert & \leq C_\lambda\lambda^n
        \label{eq:geometric_bound_of_spectral_radius_lt}
    \end{align}
    for every $n\in\N$.  The same prescribed $\lambda$ appears in the bound;
    the additional condition $\lambda<1$ is needed only when geometric decay
    is deduced downstream.  After the nonunit transfer operator has been
    isolated, this is the prescribed-rate spectral step used in
    \cite[Lemma~5.2]{Fannes1992Finitely}, but it does not establish the full
    estimate of that lemma.  The constant may depend on $\lambda$; no
    dimension-only value, in particular $C_\lambda=k^2$, is asserted.
\end{theorem}

\begin{proof}
    \leanok
    Gelfand's formula gives $N\in\N$ such that
    $\lVert a^n\rVert\leq\lambda^n$ for every $n\geq N$.  Since the spectral
    radius is nonnegative and strictly smaller than $\lambda$, one has
    $\lambda>0$.  Set
    \begin{align}
        C_\lambda
         & =1+\sum_{j=0}^{N-1}\frac{\lVert a^j\rVert}{\lambda^j}.
    \end{align}
    The eventual estimate proves the claim for $n\geq N$, while the
    corresponding summand in $C_\lambda$ proves it for $n<N$.
\end{proof}

\begin{theorem}[Pointwise geometric power bound at a prescribed rate]
    \label{thm:geometric_apply_bound_of_spectral_radius_lt}
    \lean{geometric_apply_bound_of_spectralRadius_lt}
    \leanok
    Let $V$ be a complex Banach space, let $T\colon V\to V$ be a bounded
    linear operator, and let $\lambda\geq0$.  If
    $\rho_{\mathrm{spec}}(T)<\lambda$, then there is a constant
    $C_\lambda>0$ such that
    \begin{align}
        \lVert T^n x\rVert
         & \leq C_\lambda\lambda^n\lVert x\rVert
        \label{eq:geometric_apply_bound_of_spectral_radius_lt}
    \end{align}
    for every $n\in\N$ and $x\in V$.
\end{theorem}

\begin{proof}
    \leanok
    \uses{thm:geometric_bound_of_spectral_radius_lt}
    Apply Theorem~\ref{thm:geometric_bound_of_spectral_radius_lt} in the Banach
    algebra of bounded endomorphisms of $V$, then use
    $\lVert T^n x\rVert\leq\lVert T^n\rVert\lVert x\rVert$.
\end{proof}

\begin{theorem}[FNW remainder eigenvalue gap]
    \label{thm:fnw_remainder_eigenvalue_gap}
    \lean{MPSTensor.IsPrimitiveMPS.fnwRemainder_eigenvalue_norm_lt_one}
    \leanok
    \uses{def:primitive_mps, def:fnw_transfer_map, def:fnw_limit_map}
    Let $A$ be a primitive MPS tensor of nonzero bond dimension with fixed-point
    witness $\rho$, and set
    $R_A^\rho=E_{\mathrm{FNW},A}-E_\infty^\rho$.  If $\nu$ is an eigenvalue of
    $R_A^\rho$, then
    \begin{align}
        |\nu| & <1.
        \label{eq:fnw_remainder_eigenvalue_gap}
    \end{align}
    This is the norm-independent spectral step in
    \cite[Lemma~5.2]{Fannes1992Finitely}.
\end{theorem}

\begin{proof}
    \leanok
    \uses{thm:fnw_transfer_map_trace_adjoint, thm:fnw_limit_map_trace_adjoint}
    The remainder $R_A^\rho$ is the bilinear trace adjoint of
    $\E_A-P_\rho$.  A linear map and its trace adjoint have the same
    characteristic polynomial, hence the same eigenvalues.  Primitivity places
    every eigenvalue of $\E_A-P_\rho$ strictly inside the unit disk.
\end{proof}

\begin{definition}[Trace-inverse factor]
    \label{def:fnw_trace_inverse_factor}
    \lean{MPSTensor.fnwTraceInverseFactor}
    \leanok
    Let $\rho\in\MN{D}$.  Define
    \begin{align}
        c_\rho
         & :=\operatorname{Re}\tr(\rho^{-1}).
        \label{eq:fnw_trace_inverse_factor}
    \end{align}
    This is the factor multiplying the weighted remainder norm in
    \cite[Equations~(5.9)--(5.10)]{Fannes1992Finitely}.
\end{definition}

\begin{theorem}[Positivity of the trace-inverse factor]
    \label{thm:fnw_trace_inverse_factor_pos}
    \lean{MPSTensor.fnwTraceInverseFactor_pos}
    \leanok
    \uses{def:fnw_trace_inverse_factor}
    If $\rho\in\MN{D}$ is positive definite and $D>0$, then $c_\rho>0$.
\end{theorem}

\begin{proof}
    \leanok
    \uses{def:fnw_trace_inverse_factor}
    The inverse $\rho^{-1}$ is positive definite, so its trace has strictly
    positive real part.
\end{proof}

\begin{definition}[Weighted FNW remainder data]
    \label{def:fnw_weighted_remainder_data}
    \lean{MPSTensor.fnwWeightedRemainderSpectralRadius}
    \lean{MPSTensor.fnwWeightedOperatorNorm}
    \leanok
    \uses{def:fnw_transfer_map, def:fnw_limit_map, thm:weighted-virtual-inner-product}
    Let $\rho\in\MN{D}$ be positive definite and let $\tr(\rho)\neq0$.
    Equip $\MN{D}$ with
    $\langle X,Y\rangle_\rho=\tr(\rho X^\dagger Y)$.  For a linear map
    $F\colon\MN{D}\to\MN{D}$, write
    \begin{align}
        \lVert F\rVert_{\mathrm{op},\rho}
         & :=\sup_{X\neq0}\frac{\lVert F(X)\rVert_\rho}{\lVert X\rVert_\rho}.
        \label{eq:fnw_weighted_operator_norm}
    \end{align}
    The weighted remainder spectral radius is
    \begin{align}
        r_\rho(A)
         & :=\rho_{\mathrm{spec}}\!\left(
        E_{\mathrm{FNW},A}-E_\infty^\rho\right),
        \label{eq:fnw_weighted_remainder_spectral_radius}
    \end{align}
    where the remainder is viewed as a bounded operator on this weighted
    Hilbert space.
\end{definition}

\begin{theorem}[Weighted FNW remainder spectral gap]
    \label{thm:fnw_weighted_remainder_spectral_radius_lt_one}
    \lean{MPSTensor.IsPrimitiveMPS.fnwWeightedRemainder_spectralRadius_lt_one}
    \leanok
    \uses{def:primitive_mps, def:fnw_weighted_remainder_data}
    Let $A$ be primitive with nonzero bond dimension and positive-definite
    fixed-point witness $\rho$.  Then
    \begin{align}
        r_\rho(A) & <1.
        \label{eq:fnw_weighted_remainder_spectral_radius_lt_one}
    \end{align}
\end{theorem}

\begin{proof}
    \leanok
    \uses{thm:fnw_remainder_eigenvalue_gap}
    In finite dimension the spectral radius is determined by the eigenvalues.
    Apply Theorem~\ref{thm:fnw_remainder_eigenvalue_gap} in the weighted norm.
\end{proof}

\begin{definition}[FNW mixing quantity]
    \label{def:fnw_mixing_quantity}
    \lean{MPSTensor.fnwMixingQuantity}
    \leanok
    \uses{def:fnw_trace_inverse_factor, def:fnw_weighted_remainder_data}
    Suppose that $\rho$ is positive definite and $\tr(\rho)=1$.  For
    $n\in\N$, define
    \begin{align}
        a(n)
         & :=c_\rho\,
        \bigl\lVert E_{\mathrm{FNW},A}^n-E_\infty^\rho
        \bigr\rVert_{\mathrm{op},\rho}.
        \label{eq:fnw_mixing_quantity}
    \end{align}
    This is the trace-one transfer-remainder quantity formed from the factor
    and weighted norm in
    \cite[Equations~(5.9)--(5.10)]{Fannes1992Finitely}.
\end{definition}

\begin{theorem}[Prescribed-rate bound for the FNW mixing quantity]
    \label{thm:fnw_mixing_quantity_geometric}
    \lean{MPSTensor.IsPrimitiveMPS.exists_fnwMixingQuantity_le_geometric}
    \leanok
    \uses{def:primitive_mps, def:fnw_mixing_quantity, def:fnw_weighted_remainder_data}
    Let $A$ be primitive with nonzero bond dimension and positive-definite
    trace-one fixed point $\rho$.  Let $\lambda\geq0$ satisfy
    $r_\rho(A)<\lambda$.  There is a constant $C_\lambda>0$ such that, for
    every integer $n\geq1$,
    \begin{align}
        a(n) & \leq C_\lambda\lambda^n.
        \label{eq:fnw_mixing_quantity_geometric}
    \end{align}
    This proves the prescribed-rate transfer-remainder estimate in
    \cite[Lemma~5.2]{Fannes1992Finitely}.  The condition $\lambda<1$ is needed
    only when one subsequently deduces decay to zero.
\end{theorem}

\begin{proof}
    \leanok
    \uses{thm:geometric_bound_of_spectral_radius_lt, thm:primitive_fnw_remainder_power, thm:fnw_trace_inverse_factor_pos}
    Apply Theorem~\ref{thm:geometric_bound_of_spectral_radius_lt} to the
    weighted remainder $R_A^\rho$.  This gives $C_0>0$ with
    $\lVert(R_A^\rho)^n\rVert_{\mathrm{op},\rho}
        \leq C_0\lambda^n$.  When $n\geq1$,
    Theorem~\ref{thm:primitive_fnw_remainder_power} identifies
    $(R_A^\rho)^n$ with
    $E_{\mathrm{FNW},A}^n-E_\infty^\rho$.  Multiply by $c_\rho>0$ and take
    $C_\lambda=c_\rho C_0$.
\end{proof}

The theorem does not identify $C_\lambda$ with $k^2$ or identify the FNW
auxiliary dimension $k$ with the chosen bond dimension $D$.

\subsection{The weighted FNW boundary estimate}
\label{subsec:fnw_weighted_boundary_estimate}

Let $\rho\in\MN{D}$ be positive definite with $\tr(\rho)=1$. Equip
$\MN{D}$ with the inner product
$\langle B,C\rangle_\rho=\tr(\rho B^\dagger C)$ from
\cite[Equation~(5.6)]{Fannes1992Finitely}, and let $E_{ik}$ denote the
standard matrix unit. Throughout this subsection,
$A^\mu=v(\mu)^\dagger$.

\begin{definition}[Euclidean physical-site reversal]
    \label{def:physical_site_reverse_es}
    \lean{MPSTensor.physicalSiteReverseES}
    \lean{MPSTensor.physicalSiteReverseES_apply}
    \leanok
    \uses{def:physical_site_reverse_config_equiv}
    Precomposition with the reversal $r_N$ defines a unitary operator
    \begin{align}
        \mathcal R_N\colon\C^{d^N}
         & \longrightarrow\C^{d^N}, \\
        (\mathcal R_N\psi)(\sigma)
         & =\psi(r_N(\sigma)).
        \label{eq:fnw_physical_reverse_es}
    \end{align}
\end{definition}

\begin{proof}
    \leanok
    \uses{def:physical_site_reverse_config_equiv}
    A permutation of the standard coordinates preserves the Euclidean inner
    product, and evaluation gives
    (\ref{eq:fnw_physical_reverse_es}).
\end{proof}

\begin{definition}[Weighted FNW boundary operator]
    \label{def:fnw_boundary_clm}
    \lean{MPSTensor.fnwBoundaryMapCLM}
    \lean{MPSTensor.fnwBoundaryMapCLM_apply}
    \lean{MPSTensor.fnwBoundaryMapCLM_eq_physicalSiteReverseES_comp}
    \leanok
    \uses{def:fnw_boundary_map, def:physical_site_reverse_es,
        thm:weighted-virtual-inner-product}
    The boundary map of \cite[Equation~(5.5)]{Fannes1992Finitely} defines a
    bounded linear operator
    \begin{align}
        F_N\colon(\MN{D},\langle\cdot,\cdot\rangle_\rho)
         & \longrightarrow\C^{d^N},                \\
        F_N(B)(\sigma)
         & =\tr\!\left(B(v^\sigma)^\dagger\right).
        \label{eq:fnw_boundary_clm_apply}
    \end{align}
    In the ground-space convention,
    \begin{align}
        F_N & =\mathcal R_N\circ\Gamma_N^A.
        \label{eq:fnw_boundary_clm_ground_space}
    \end{align}
\end{definition}

\begin{proof}
    \leanok
    \uses{thm:fnw_boundary_map_eq_physical_site_reverse_comp_ground_space_map,
        def:physical_site_reverse_es}
    The algebraic identity
    $F_N^v=R_N\circ\Gamma_N^A$ from
    Theorem~\ref{thm:fnw_boundary_map_eq_physical_site_reverse_comp_ground_space_map}
    is unchanged after equipping the domain with the weighted norm and the
    codomain with the Euclidean norm. The reversal becomes the unitary
    operator $\mathcal R_N$.
\end{proof}

\begin{theorem}[FNW boundary scalar product]
    \label{thm:fnw_boundary_inner}
    \lean{MPSTensor.inner_fnwBoundaryMapCLM}
    \lean{MPSTensor.inner_fnwBoundaryMapCLM_eq_sum_fnwTransferMap}
    \leanok
    \uses{def:fnw_boundary_clm, def:fnw_transfer_map}
    For all $B,C\in\MN{D}$,
    \begin{align}
        \langle F_N(B),F_N(C)\rangle
         & =\sum_{i,k}
        \left(B^\dagger E_{\mathrm{FNW},A}^N(E_{ik})C\right)_{ik}
        \label{eq:fnw_boundary_inner_choi} \\
         & =\sum_{i,k}
        \left(E_{\mathrm{FNW},A}^N
        \left(B^\dagger E_{ik}C\right)\right)_{ik}.
        \label{eq:fnw_boundary_inner_source}
    \end{align}
    These are the two forms of the scalar-product identity in
    \cite[Equation~(5.8)]{Fannes1992Finitely}. The second equality preserves
    the source orientation $B^\dagger E_{ik}C$.
\end{theorem}

\begin{proof}
    \leanok
    \uses{thm:fnw_boundary_map_apply, def:fnw_transfer_map}
    Expand both boundary coefficients using
    (\ref{eq:fnw_boundary_clm_apply}), interchange the finite sums, and write
    the two traces as sums of diagonal entries. Multiplication by the matrix
    unit $E_{ik}$ gives respectively
    (\ref{eq:fnw_boundary_inner_choi}) and
    (\ref{eq:fnw_boundary_inner_source}).
\end{proof}

\begin{theorem}[Exact weighted boundary defect]
    \label{thm:fnw_boundary_defect}
    \lean{MPSTensor.inner_fnwBoundaryMapCLM_sub_rhoWeighted}
    \leanok
    \uses{thm:fnw_boundary_inner, def:fnw_limit_map,
        thm:weighted-virtual-inner-product}
    For all $B,C\in\MN{D}$,
    \begin{align}
        \langle F_N(B),F_N(C)\rangle-\langle B,C\rangle_\rho
         & =\sum_{i,k}
        \left(\left(E_{\mathrm{FNW},A}^N-E_\infty^\rho\right)
        \left(B^\dagger E_{ik}C\right)\right)_{ik}.
        \label{eq:fnw_boundary_defect}
    \end{align}
\end{theorem}

\begin{proof}
    \leanok
    \uses{thm:fnw_boundary_inner, def:fnw_limit_map,
        thm:weighted-virtual-inner-product}
    The trace-one limiting map satisfies
    \begin{align}
        \sum_{i,k}
        \left(E_\infty^\rho(B^\dagger E_{ik}C)\right)_{ik}
         & =\tr(\rho B^\dagger C)=\langle B,C\rangle_\rho.
    \end{align}
    Subtract this identity from
    (\ref{eq:fnw_boundary_inner_source}).
\end{proof}

\begin{theorem}[Weighted FNW boundary near-isometry]
    \label{thm:fnw_boundary_near_isometry_exact}
    \lean{MPSTensor.norm_inner_fnwBoundaryMapCLM_sub_rhoWeighted_le}
    \leanok
    \uses{thm:fnw_boundary_defect, def:fnw_weighted_remainder_data}
    For all $B,C\in\MN{D}$,
    \begin{align}
        \bigl|\langle F_N(B),F_N(C)\rangle
        -\langle B,C\rangle_\rho\bigr|
         & \leq\operatorname{Re}\tr(\rho^{-1})
        \bigl\lVert E_{\mathrm{FNW},A}^N-E_\infty^\rho
        \bigr\rVert_{\mathrm{op},\rho}
        \lVert B\rVert_\rho\lVert C\rVert_\rho.
        \label{eq:fnw_boundary_near_isometry_exact}
    \end{align}
    This is the estimate in
    \cite[Equation~(5.9)]{Fannes1992Finitely}. It contains no factor depending
    only on the dimension.
\end{theorem}

\begin{proof}
    \leanok
    \uses{thm:fnw_boundary_defect, def:fnw_weighted_remainder_data}
    Regard the sum in (\ref{eq:fnw_boundary_defect}) as the trace of the
    endomorphism obtained by composing the transfer remainder with
    $X\mapsto B^\dagger XC$. Weighted trace duality bounds this trace by
    $\operatorname{Re}\tr(\rho^{-1})$ times the weighted operator norm of the
    remainder and the two weighted boundary norms.
\end{proof}

\begin{theorem}[Right-boundary successor recursion]
    \label{thm:fnw_boundary_norm_sq_succ}
    \lean{MPSTensor.fnwBoundaryMapCLM_norm_sq_succ}
    \leanok
    \uses{def:fnw_boundary_clm}
    For every $B\in\MN{D}$,
    \begin{align}
        \lVert F_{N+1}(B)\rVert^2
         & =\sum_{\mu=0}^{d-1}\lVert F_N(BA^\mu)\rVert^2.
        \label{eq:fnw_boundary_norm_sq_succ}
    \end{align}
    This recursion follows from the boundary formula in
    \cite[Equation~(5.5)]{Fannes1992Finitely} and is used in the proof of
    \cite[Lemma~5.3]{Fannes1992Finitely}. Here
    $A^\mu=v(\mu)^\dagger$.
\end{theorem}

\begin{proof}
    \leanok
    \uses{thm:fnw_boundary_map_apply, def:eval_word}
    Split a configuration of length $N+1$ into its first $N$ entries and its
    final entry $\mu$. The adjointed word identity gives
    $F_{N+1}(B)(\sigma,\mu)=F_N(BA^\mu)(\sigma)$. Summing the squared absolute
    values over $\sigma$ and $\mu$ proves
    (\ref{eq:fnw_boundary_norm_sq_succ}).
\end{proof}

The preceding results establish Equations~(5.8)--(5.9) and the right-boundary
recursion, but not Lemmas~5.3 or~6.2 of
Ref.~\cite{Fannes1992Finitely}. In the source proof of Lemma~6.2,
Equation~(6.5) contributes $a(m)/a_-(m)$, while the subsequent orthogonality
estimates contribute $a(m)^2/a_-(m)$. Their sum is
$a(m)(1+a(m))/a_-(m)$.

\begin{theorem}[Geometric power bound below unit spectral radius]
    \label{thm:geometric_bound_of_spectral_radius_lt_one}
    \lean{geometric_bound_of_spectralRadius_lt_one}
    \leanok
    Let $V$ be a complex Banach space and let $T\colon V\to V$ be a bounded
    linear operator.  If $\rho_{\mathrm{spec}}(T)<1$, then there are real
    constants $C,r$ with $C>0$ and $0<r<1$ such that
    \begin{align}
        \lVert T^n\rVert & \leq C r^n
        \label{eq:geometric_bound_of_spectral_radius_lt_one}
    \end{align}
    for every $n\in\N$.
\end{theorem}

\begin{proof}
    \leanok
    \uses{thm:geometric_bound_of_spectral_radius_lt}
    Choose $r$ strictly between the spectral radius and one, then apply
    Theorem~\ref{thm:geometric_bound_of_spectral_radius_lt} with
    $\lambda=r$.
\end{proof}

\begin{theorem}[Pointwise geometric power bound below unit spectral radius]
    \label{thm:geometric_apply_bound_of_spectral_radius_lt_one}
    \lean{geometric_apply_bound_of_spectralRadius_lt_one}
    \leanok
    \uses{thm:geometric_bound_of_spectral_radius_lt_one}
    Let $V$ be a complex Banach space and let $T\colon V\to V$ be a bounded
    linear operator.  If $\rho_{\mathrm{spec}}(T)<1$, then there are real
    constants $C,r$ with $C>0$ and $0<r<1$ such that
    \begin{align}
        \lVert T^n x\rVert & \leq C r^n\lVert x\rVert
        \label{eq:geometric_apply_bound_of_spectral_radius_lt_one}
    \end{align}
    for every $n\in\N$ and $x\in V$.
\end{theorem}

\begin{proof}
    \leanok
    \uses{thm:geometric_bound_of_spectral_radius_lt_one}
    Apply Theorem~\ref{thm:geometric_bound_of_spectral_radius_lt_one} and the
    defining operator-norm inequality to $T^n x$.
\end{proof}

\begin{theorem}[Transfer--Gram Choi reshuffling]
    \label{thm:ground_space_gram_choi_reshuffling}
    \lean{MPSTensor.inner_single_groundSpaceGram_single_eq_rectangularChoi}
    \leanok
    \uses{def:ground_space_gram, def:transfer_map}
    Let $\mathbf e_{(j,i)}:=\operatorname{single}_{(j,i)}(1)$ denote the
    standard unit vector in $\ell^2(\Fin{D}\times\Fin{D})$, and let
    $\mathcal E_A$ be the transfer map.  Then
    \begin{align}
        \bigl\langle \mathbf e_{(b,a)},
        \mathcal K_L(\mathbf e_{(e,c)})\bigr\rangle_{\ell^2}
         & =J(\mathcal E_A^L)_{(c,e),(a,b)}.
    \end{align}
    Under Frobenius vectorization, $\mathbf e_{(b,a)}=U(E_{ab})$ and
    $\mathbf e_{(e,c)}=U(E_{ce})$, explaining the reversed coordinate order.
    Thus the Gram operator is the indicated Choi reshuffling; in particular it
    is not obtained by replacing the right-hand side with
    $\langle E_{ab},\mathcal E_A^L(E_{ce})\rangle_{\mathrm{HS}}$.
\end{theorem}

\begin{proof}
    \leanok
    \uses{thm:ground_space_map_es_single, lem:ground_space_map_apply, thm:mixed_pow}
    Move the adjoint in
    $\mathcal K_L=(\Gamma_L^{\mathrm{ES}})^\dagger\Gamma_L^{\mathrm{ES}}$
    across the
    inner product and expand the two boundary-map coefficients.  For each
    physical word $\sigma$, the resulting summand is
    $(A^\sigma)_{ec}\,\overline{(A^\sigma)_{ba}}$.  On the other hand,
    expanding the transfer power on $E_{ca}$ gives
    \begin{align}
        \bigl(\mathcal E_A^L(E_{ca})\bigr)_{eb}
         & =\sum_\sigma (A^\sigma)_{ec}\,
        \overline{(A^\sigma)_{ba}}.
    \end{align}
    The defining Choi indexing
    $J(\mathcal L)_{(c,e),(a,b)}=(\mathcal L(E_{ca}))_{eb}$ now yields the
    stated reshuffling.
\end{proof}

\begin{theorem}[Ground-space Gram operator as a reshuffled Choi matrix]
    \label{thm:ground_space_gram_eq_gram_reshuffle}
    \lean{MPSTensor.groundSpaceGram_eq_gramReshuffle}
    \leanok
    \uses{def:ground_space_gram, def:gram_reshuffle, def:transfer_map}
    The ground-space Gram operator is exactly the operator-level reshuffling of
    the Choi matrix of the iterated transfer map:
    \begin{align}
        \mathcal K_L
         & =\mathscr R\!\left(\mathcal E_A^L\right).
        \label{eq:ground_space_gram_eq_gram_reshuffle}
    \end{align}
    The equality is an equality of operators on
    $\ell^2(\Fin{D}\times\Fin{D})$.  Its entries are those in
    (\ref{eq:gram_reshuffle_single}), rather than Hilbert--Schmidt pairings with
    $\mathcal E_A^L$.
\end{theorem}

\begin{proof}
    \leanok
    \uses{thm:ground_space_gram_choi_reshuffling, thm:gram_reshuffle_single}
    The two operators have equal matrix coefficients on every pair of standard
    basis vectors.  Hence they are equal.
\end{proof}

\begin{theorem}[Geometric convergence of the ground-space Gram operator]
    \label{thm:ground_space_gram_geometric_fixed_point}
    \lean{MPSTensor.groundSpaceGram_sub_fixedPointProj_norm_sq_le_geometric}
    \leanok
    \uses{def:ground_space_gram, def:gram_reshuffle, def:transfer_map}
    Let $A$ be an MPS tensor of bond dimension $D$, let $\rho\in\MN{D}$ satisfy
    $\tr(\rho)\neq0$, and let $P_\rho$ be the fixed-point projection
    \begin{align}
        P_\rho(X) & =\frac{\tr(X)}{\tr(\rho)}\rho.
    \end{align}
    Suppose that $\mathcal E_A$ is trace-preserving,
    $\mathcal E_A(\rho)=\rho$, and
    \begin{align}
        \rho_{\mathrm{spec}}(\mathcal E_A-P_\rho) & <1.
    \end{align}
    Then there are real constants $C,r$ with $C>0$ and $0<r<1$ such that,
    for every $n\geq1$,
    \begin{align}
        \bigl\lVert\mathcal K_n-\mathscr R(P_\rho)\bigr\rVert^2
         & \leq D^3\bigl(Cr^n\bigr)^2.
        \label{eq:ground_space_gram_geometric_fixed_point}
    \end{align}
    No positivity, injectivity, or uniqueness hypothesis is imposed.
\end{theorem}

\begin{proof}
    \leanok
    \uses{thm:geometric_bound_of_spectral_radius_lt_one, thm:ground_space_gram_eq_gram_reshuffle, thm:gram_reshuffle_sub, thm:gram_reshuffle_norm_sq_le_card_cube_mul_operator_norm_sq}
    Write $N=\mathcal E_A-P_\rho$.  The power decomposition gives
    $\mathcal E_A^n=P_\rho+N^n$ for $n\geq1$, while
    Theorem~\ref{thm:geometric_bound_of_spectral_radius_lt_one} gives
    $\lVert N^n\rVert\leq Cr^n$.  Use the exact Gram reshuffling identity and
    Theorem~\ref{thm:gram_reshuffle_sub} to rewrite the left-hand side as
    $\lVert\mathscr R(N^n)\rVert^2$, then apply
    Theorem~\ref{thm:gram_reshuffle_norm_sq_le_card_cube_mul_operator_norm_sq}.
\end{proof}

\begin{theorem}[Primitive-MPS geometric Gram convergence]
    \label{thm:primitive_ground_space_gram_geometric_fixed_point}
    \lean{MPSTensor.IsPrimitiveMPS.groundSpaceGram_sub_fixedPointProj_norm_sq_le_geometric}
    \leanok
    \uses{def:primitive_mps, thm:ground_space_gram_geometric_fixed_point}
    Let $A$ be a primitive MPS tensor with fixed-point witness $\rho$. Then there
    are real constants $C,r$ with $C>0$ and $0<r<1$ such that, for every
    $n\geq1$,
    \begin{align}
        \bigl\lVert\mathcal K_n-\mathscr R(P_\rho)\bigr\rVert^2
         & \leq D^3\bigl(Cr^n\bigr)^2.
    \end{align}
\end{theorem}

\begin{proof}
    \leanok
    \uses{thm:ground_space_gram_geometric_fixed_point}
    Positive semidefiniteness and nonvanishing of $\rho$ imply
    $\tr(\rho)\neq0$. The trace-preserving identity, fixed-point equation, and
    complementary spectral-radius bound are the remaining fields of the
    primitive-MPS hypothesis, so
    Theorem~\ref{thm:ground_space_gram_geometric_fixed_point} applies.
\end{proof}

\begin{theorem}[Ground-space Gram convergence]
    \label{thm:ground_space_gram_tendsto_fixed_point}
    \lean{MPSTensor.groundSpaceGram_tendsto_gramReshuffle_fixedPointProj}
    \leanok
    \uses{def:ground_space_gram, def:gram_reshuffle, def:transfer_map}
    Let $A$ be an MPS tensor of bond dimension $D$, let $\rho\in\MN{D}$ satisfy
    $\tr(\rho)\neq0$, and let $P_\rho$ be the corresponding fixed-point
    projection.  If $\mathcal E_A$ is trace-preserving,
    $\mathcal E_A(\rho)=\rho$, and
    $\rho_{\mathrm{spec}}(\mathcal E_A-P_\rho)<1$, then
    \begin{align}
        \mathcal K_n & \longrightarrow\mathscr R(P_\rho)
        \label{eq:ground_space_gram_tendsto_fixed_point}
    \end{align}
    in operator norm as $n\to\infty$.
\end{theorem}

\begin{proof}
    \leanok
    \uses{thm:ground_space_gram_geometric_fixed_point}
    The right-hand side of
    (\ref{eq:ground_space_gram_geometric_fixed_point}) tends to zero because
    $0<r<1$.  Taking square roots gives the asserted norm convergence.
\end{proof}

\begin{theorem}[Primitive-MPS ground-space Gram convergence]
    \label{thm:primitive_ground_space_gram_tendsto_fixed_point}
    \lean{MPSTensor.IsPrimitiveMPS.groundSpaceGram_tendsto_gramReshuffle_fixedPointProj}
    \leanok
    \uses{def:primitive_mps, thm:ground_space_gram_tendsto_fixed_point}
    If $A$ is a primitive MPS tensor with fixed-point witness $\rho$, then
    \begin{align}
        \mathcal K_n & \longrightarrow\mathscr R(P_\rho)
    \end{align}
    in operator norm as $n\to\infty$.
\end{theorem}

\begin{proof}
    \leanok
    \uses{thm:ground_space_gram_tendsto_fixed_point}
    The primitive-MPS data supply the nonzero trace, trace preservation,
    fixed-point equation, and complementary spectral-radius hypothesis of
    Theorem~\ref{thm:ground_space_gram_tendsto_fixed_point}.
\end{proof}

The limit in Theorem~\ref{thm:ground_space_gram_tendsto_fixed_point} is
the reshuffled fixed-point projection $\mathscr R(P_\rho)$, not the
identity in unweighted Frobenius coordinates.  The limiting metric is
generally nonidentity.
Under a positive-definite fixed-point hypothesis, the inverse estimates
below perturb around $\mathscr R(P_\rho)$ rather than around the identity.
They do not yet compare the physical range projectors on overlapping
windows.

\begin{theorem}[Exact limiting Gram metric]
    \label{thm:fixed_point_gram_metric_exact}
    \lean{Matrix.gramReshuffleMatrix_fixedPointProj_apply}
    \lean{Matrix.gramReshuffleMatrix_fixedPointProj}
    \lean{Matrix.gramReshuffle_fixedPointProj_frobeniusEquivEuclidean_apply}
    \leanok
    \uses{def:gram_reshuffle}
    Let $\rho\in\MN{D}$ have nonzero trace and define
    $P_\rho(X):=(\tr X/\tr\rho)\rho$.  Then
    \begin{align}
        R(P_\rho)
         & =(\tr\rho)^{-1}(\rho^{\mathsf T}\otimes I_D),
        \label{eq:fixed_point_gram_metric_matrix}        \\
        \mathscr R(P_\rho)U(X)
         & =U\!\left((\tr\rho)^{-1}X\rho\right).
        \label{eq:fixed_point_gram_metric_action}
    \end{align}
\end{theorem}

\begin{proof}
    \leanok
    \uses{def:gram_reshuffle}
    For matrix units $E_{ca}$, direct evaluation gives
    \begin{align}
        R(P_\rho)_{(b,a),(e,c)}
         & =(\tr\rho)^{-1}\rho_{eb}\,\delta_{ca}.
    \end{align}
    These are exactly the entries in
    (\ref{eq:fixed_point_gram_metric_matrix}).  Applying that Kronecker product
    to the column--row vectorization of $X$ gives
    (\ref{eq:fixed_point_gram_metric_action}).
\end{proof}

\begin{theorem}[Invertibility of the limiting Gram metric]
    \label{thm:fixed_point_gram_metric_is_unit}
    \lean{Matrix.gramReshuffleMatrix_fixedPointProj_posDef}
    \lean{Matrix.gramReshuffle_fixedPointProj_isUnit}
    \lean{Matrix.gramReshuffle_fixedPointProj_injective}
    \leanok
    \uses{def:gram_reshuffle}
    If $D>0$ and $\rho\in\MN{D}$ is positive definite, then
    $R(P_\rho)$ is positive definite and $\mathscr R(P_\rho)$ is invertible,
    hence injective.
\end{theorem}

\begin{proof}
    \leanok
    \uses{def:gram_reshuffle, thm:fixed_point_gram_metric_exact}
    Positive definiteness gives $\tr\rho>0$.  By
    (\ref{eq:fixed_point_gram_metric_matrix}),
    \begin{align}
        R(P_\rho)
         & =(\tr\rho)^{-1}(\rho^{\mathsf T}\otimes I_D)>0.
    \end{align}
    Hence its associated operator $\mathscr R(P_\rho)$ is invertible.
\end{proof}

\begin{theorem}[Arbitrary-base Neumann inverse bound]
    \label{thm:arbitrary_base_neumann_inverse_bound}
    \lean{NormedRing.isUnit_and_norm_inverse_le_of_norm_inverse_mul_sub_le}
    \lean{NormedRing.isUnit_and_norm_inverse_le_of_norm_mul_norm_sub_le}
    \leanok
    Let $K_0$ be invertible in a normed ring with convergent geometric series
    and $\lVert I\rVert=1$.  If $a<1$ and
    \begin{align}
        \lVert K_0^{-1}(K-K_0)\rVert & \le a,
        \label{eq:arbitrary_base_relative_perturbation}
    \end{align}
    then $K$ is invertible and
    \begin{align}
        \lVert K^{-1}\rVert
         & \le(1-a)^{-1}\lVert K_0^{-1}\rVert.
        \label{eq:arbitrary_base_inverse_bound}
    \end{align}
    The same conclusion follows from
    $\lVert K_0^{-1}\rVert\,\lVert K-K_0\rVert\le a$.
\end{theorem}

\begin{proof}
    \leanok
    \uses{thm:neumann_inverse_norm_bound}
    Put $E:=K_0^{-1}(K-K_0)$.  Then
    $K=K_0(I+E)$ and $\lVert E\rVert<1$, so
    \begin{align}
        K^{-1} & =(I+E)^{-1}K_0^{-1},                             \\
        \lVert K^{-1}\rVert
               & \le\lVert(I+E)^{-1}\rVert\,\lVert K_0^{-1}\rVert
        \le(1-a)^{-1}\lVert K_0^{-1}\rVert.
    \end{align}
    Submultiplicativity reduces the product-form hypothesis to
    (\ref{eq:arbitrary_base_relative_perturbation}).
\end{proof}

\begin{theorem}[Arbitrary-base inverse-displacement bound]
    \label{thm:arbitrary_base_neumann_inverse_displacement}
    \lean{NormedRing.norm_inverse_sub_inverse_le_of_norm_inverse_mul_sub_le}
    \lean{NormedRing.norm_inverse_sub_inverse_le_of_norm_mul_norm_sub_le}
    \leanok
    Let $K_0$ be invertible in a normed ring with convergent geometric series
    and $\lVert I\rVert=1$.  If $a<1$ and
    (\ref{eq:arbitrary_base_relative_perturbation}) holds, then
    \begin{align}
        \lVert K^{-1}-K_0^{-1}\rVert
         & \le(1-a)^{-1}a\,\lVert K_0^{-1}\rVert.
        \label{eq:arbitrary_base_inverse_displacement}
    \end{align}
    The same estimate follows from
    $\lVert K_0^{-1}\rVert\,\lVert K-K_0\rVert\le a$.
\end{theorem}

\begin{proof}
    \leanok
    \uses{thm:neumann_inverse_norm_bound}
    With $E:=K_0^{-1}(K-K_0)$, the resolvent identity is
    \begin{align}
        K^{-1}-K_0^{-1}
         & =-(I+E)^{-1}E K_0^{-1}.
    \end{align}
    Since $\lVert(I+E)^{-1}\rVert\le(1-a)^{-1}$ and
    $\lVert E\rVert\le a$, taking norms gives
    (\ref{eq:arbitrary_base_inverse_displacement}).  Submultiplicativity again
    gives the product-form variant.
\end{proof}

\begin{theorem}[Continuity of inversion at an invertible limit]
    \label{thm:inverse_tendsto_at_unit}
    \lean{NormedRing.inverse_tendsto_of_tendsto_of_isUnit}
    \leanok
    Let $K_i$ be a net in a normed ring with convergent geometric series.  If
    $K_i\to K_0$ and $K_0$ is invertible, then
    \begin{align}
        K_i^{-1} & \longrightarrow K_0^{-1}.
    \end{align}
\end{theorem}

\begin{proof}
    \leanok
    Write $K_0$ as the image of a unit.  Continuity of inversion on the unit
    group, followed by the convergence $K_i\to K_0$, gives the result.
\end{proof}

\begin{definition}[Inverse Gram operator]
    \label{def:inverse_gram}
    \lean{ContinuousLinearMap.inverseGram}
    \leanok
    For an injective continuous linear map $T$ with finite-dimensional domain,
    let $S_T$ be the continuous endomorphism induced by the inverse linear
    equivalence of the Gram endomorphism $T^\dagger T$.
\end{definition}

\begin{theorem}[Inverse-Gram identities]
    \label{thm:inverse_gram_identities}
    \lean{ContinuousLinearMap.inverseGram_comp_adjoint_comp_self}
    \lean{ContinuousLinearMap.adjoint_comp_self_comp_inverseGram}
    \lean{ContinuousLinearMap.inverseGram_eq_inverse}
    \lean{ContinuousLinearMap.inverseGram_eq_ringInverse}
    \lean{ContinuousLinearMap.norm_T_comp_inverseGram_eq_sqrt}
    \lean{ContinuousLinearMap.norm_inverseGram_comp_adjoint_eq_sqrt}
    \leanok
    \uses{def:inverse_gram}
    Let $\mathbb K$ be a real or complex scalar field, let $E$ be a
    finite-dimensional complete $\mathbb K$-Hilbert space, and let $F$ be a
    complete $\mathbb K$-Hilbert space.  If $T\colon E\to F$ is an injective
    continuous linear map, then the inverse Gram operator is a two-sided
    inverse of the Gram endomorphism:
    \begin{align}
        S_T(T^\dagger T) & =I_E,
        \label{eq:inverse_gram_left} \\
        (T^\dagger T)S_T & =I_E.
        \label{eq:inverse_gram_right}
    \end{align}
    Consequently, $S_T$ is both the canonical inverse and the ring inverse of
    $T^\dagger T$.  The two pseudoinverse factors satisfy
    \begin{align}
        \lVert TS_T\rVert         & =\sqrt{\lVert S_T\rVert}, \\
        \lVert S_TT^\dagger\rVert & =\sqrt{\lVert S_T\rVert}.
    \end{align}
\end{theorem}

\begin{proof}
    \leanok
    \uses{def:inverse_gram}
    By definition, $S_T$ is induced by the inverse linear equivalence of the
    injective endomorphism $T^\dagger T$.  The inverse-equivalence identities
    give (\ref{eq:inverse_gram_left}) and
    (\ref{eq:inverse_gram_right}).  Together these equations characterize the
    canonical inverse of a continuous linear endomorphism, which agrees with
    its inverse in the endomorphism ring.  For the norm identities, use the
    $C^*$-identity and the two inverse equations to compute the squares of
    $\lVert TS_T\rVert$ and $\lVert S_TT^\dagger\rVert$.
\end{proof}

The next three results give quantitative inverse bounds for the nonidentity
limiting metric $\mathscr R(P_\rho)$.  They require the fixed-point witness
$\rho$ to be positive definite and preserve the squared-bound normalization of
Theorem~\ref{thm:primitive_ground_space_gram_geometric_fixed_point}.  They do
not estimate a contraction between physical projectors on overlapping
windows.

\begin{theorem}[Pointwise geometric Gram inverse bounds]
    \label{thm:primitive_ground_space_gram_geometric_inverse_bounds}
    \lean{MPSTensor.IsPrimitiveMPS.groundSpaceGram_geometric_inverse_bounds}
    \leanok
    \uses{def:primitive_mps}
    Let $A$ be a primitive MPS tensor of nonzero bond dimension whose
    fixed-point witness $\rho$ is positive definite.  Put
    \begin{align}
        \mathcal K_\infty & =\mathscr R(P_\rho),     &
        \mathcal I_\infty & =\mathcal K_\infty^{-1}.
    \end{align}
    There are constants $g,c,r>0$ with $r<1$ and
    $c=\lVert\mathcal I_\infty\rVert g$ such that, for every $n\geq1$,
    \begin{align}
        \lVert\mathcal K_n-\mathcal K_\infty\rVert
         & \leq gr^n.
        \label{eq:primitive_ground_space_gram_geometric_bound}
    \end{align}
    Whenever $cr^n<1$, the operator $\mathcal K_n$ is invertible and
    \begin{align}
        \lVert\mathcal K_n^{-1}\rVert
         & \leq\frac{\lVert\mathcal I_\infty\rVert}{1-cr^n},
        \label{eq:primitive_ground_space_gram_geometric_inverse_norm} \\
        \lVert\mathcal K_n^{-1}-\mathcal I_\infty\rVert
         & \leq\frac{cr^n}{1-cr^n}\,
        \lVert\mathcal I_\infty\rVert.
        \label{eq:primitive_ground_space_gram_geometric_inverse_displacement}
    \end{align}
\end{theorem}

\begin{proof}
    \leanok
    \uses{thm:primitive_ground_space_gram_geometric_fixed_point, thm:fixed_point_gram_metric_is_unit, thm:arbitrary_base_neumann_inverse_bound, thm:arbitrary_base_neumann_inverse_displacement}
    Start from the squared estimate
    $\lVert\mathcal K_n-\mathcal K_\infty\rVert^2
        \leq D^3(Cr^n)^2$.
    Taking square roots gives
    $g=\sqrt{D^3}\,C$ without any additional dimension loss, and hence
    (\ref{eq:primitive_ground_space_gram_geometric_bound}).  Positive
    definiteness makes $\mathcal K_\infty$ invertible.  The relative error is
    at most $\lVert\mathcal I_\infty\rVert gr^n=cr^n$, so the two
    arbitrary-base Neumann estimates give
    (\ref{eq:primitive_ground_space_gram_geometric_inverse_norm}) and
    (\ref{eq:primitive_ground_space_gram_geometric_inverse_displacement}).
\end{proof}

\begin{theorem}[Pointwise geometric inverse-Gram bounds]
    \label{thm:primitive_ground_space_map_geometric_inverse_gram_bounds}
    \lean{MPSTensor.IsPrimitiveMPS.groundSpaceMapES_geometric_inverseGram_bounds}
    \leanok
    \uses{def:ground_space_map_es, def:inverse_gram, thm:primitive_ground_space_gram_geometric_inverse_bounds, thm:inverse_gram_identities}
    Under the hypotheses of
    Theorem~\ref{thm:primitive_ground_space_gram_geometric_inverse_bounds},
    there are constants $g,c,r>0$ with $r<1$ and
    $c=\lVert\mathcal I_\infty\rVert g$ such that
    (\ref{eq:primitive_ground_space_gram_geometric_bound}) holds for every
    $n\geq1$.  If $cr^n<1$, then the physical boundary map
    $\Gamma_n^{\mathrm{ES}}$ is injective, its inverse Gram operator equals
    $\mathcal K_n^{-1}$, and
    \begin{align}
        \lVert S_{\Gamma_n^{\mathrm{ES}}}\rVert
         & \leq\frac{\lVert\mathcal I_\infty\rVert}{1-cr^n}, \\
        \lVert S_{\Gamma_n^{\mathrm{ES}}}-\mathcal I_\infty\rVert
         & \leq\frac{cr^n}{1-cr^n}\,
        \lVert\mathcal I_\infty\rVert.
    \end{align}
\end{theorem}

\begin{proof}
    \leanok
    \uses{thm:primitive_ground_space_gram_geometric_inverse_bounds, thm:inverse_gram_identities}
    Invertibility of
    $(\Gamma_n^{\mathrm{ES}})^\dagger\Gamma_n^{\mathrm{ES}}=\mathcal K_n$
    implies injectivity of $\Gamma_n^{\mathrm{ES}}$.  The inverse-Gram
    identities identify $S_{\Gamma_n^{\mathrm{ES}}}$ with
    $\mathcal K_n^{-1}$, so the preceding pointwise bounds apply unchanged.
\end{proof}

\begin{theorem}[Geometric smallness at the C3 correction lengths]
    \label{thm:geometric_smallness_at_c3_lengths}
    \lean{MPSTensor.geometric_smallness_at_c3_lengths}
    \leanok
    Let $c>0$, let $0<r<1$, and let $l,K\in\mathbb N$.  If $cr^l<1$,
    then
    \begin{align}
        cr^{l+1}   & <1, \\
        cr^{K+l}   & <1, \\
        cr^{K+l+1} & <1.
    \end{align}
\end{theorem}

\begin{proof}
    \leanok
    Since $0<r<1$, the sequence $r^m$ is nonincreasing.  Each of
    $l+1$, $K+l$, and $K+l+1$ is at least $l$, so every displayed inequality
    follows from $cr^l<1$.
\end{proof}

\begin{theorem}[Eventual inverse-Gram control in the limiting metric]
    \label{thm:primitive_ground_space_gram_eventual_inverse_bound}
    \lean{MPSTensor.IsPrimitiveMPS.eventually_groundSpaceGram_isUnit_and_inverse_bound}
    \leanok
    \uses{def:primitive_mps, thm:primitive_ground_space_gram_tendsto_fixed_point}
    Let $A$ be a primitive MPS tensor whose fixed-point witness $\rho$ is
    positive definite, and put
    $\mathcal K_\infty=\mathscr R(P_\rho)$.  For every $0<a<1$, all
    sufficiently large $n$ satisfy
    \begin{align}
        \mathcal K_n & \in\operatorname{GL},
        \notag                               \\
        \bigl\lVert\mathcal K_n^{-1}-\mathcal K_\infty^{-1}\bigr\rVert
                     & \leq
        \frac{a}{1-a}\,\bigl\lVert\mathcal K_\infty^{-1}\bigr\rVert.
        \label{eq:primitive_ground_space_gram_eventual_inverse_bound}
    \end{align}
\end{theorem}

\begin{proof}
    \leanok
    \uses{thm:primitive_ground_space_gram_tendsto_fixed_point, thm:fixed_point_gram_metric_is_unit, thm:arbitrary_base_neumann_inverse_bound, thm:arbitrary_base_neumann_inverse_displacement}
    The positive-definite fixed point makes
    $\mathcal K_\infty$ invertible.  Gram convergence implies that eventually
    \begin{align}
        \bigl\lVert\mathcal K_\infty^{-1}\bigr\rVert
        \bigl\lVert\mathcal K_n-\mathcal K_\infty\bigr\rVert & \leq a.
    \end{align}
    Factoring $\mathcal K_n$ around the invertible base
    $\mathcal K_\infty$ and applying the arbitrary-base Neumann estimate gives
    (\ref{eq:primitive_ground_space_gram_eventual_inverse_bound}).
\end{proof}

\begin{theorem}[Convergence of finite-volume inverse Gram operators]
    \label{thm:primitive_ground_space_gram_inverse_tendsto}
    \lean{MPSTensor.IsPrimitiveMPS.groundSpaceGram_ringInverse_tendsto}
    \leanok
    \uses{def:primitive_mps, thm:primitive_ground_space_gram_tendsto_fixed_point}
    Under the same primitive-MPS and positive-definite fixed-point hypotheses,
    \begin{align}
        \mathcal K_n^{-1} & \longrightarrow\mathcal K_\infty^{-1}
        \label{eq:primitive_ground_space_gram_inverse_tendsto}
    \end{align}
    in operator norm.
\end{theorem}

\begin{proof}
    \leanok
    \uses{thm:primitive_ground_space_gram_tendsto_fixed_point, thm:fixed_point_gram_metric_is_unit, thm:inverse_tendsto_at_unit}
    Inversion is continuous at the invertible operator
    $\mathcal K_\infty$, so the claim follows from
    $\mathcal K_n\to\mathcal K_\infty$.
\end{proof}

\begin{theorem}[Eventual injectivity and inverse-Gram control]
    \label{thm:primitive_ground_space_map_eventual_inverse_gram_bound}
    \lean{MPSTensor.IsPrimitiveMPS.eventually_groundSpaceMapES_injective_and_inverseGram_bound}
    \leanok
    \uses{def:ground_space_map_es, def:inverse_gram, thm:primitive_ground_space_gram_eventual_inverse_bound, thm:inverse_gram_identities}
    Under the same hypotheses, for every $0<a<1$ and all sufficiently large
    $n$, the physical boundary map $\Gamma_n^{\mathrm{ES}}$ is injective and
    its inverse Gram operator is $\mathcal K_n^{-1}$.  Moreover,
    \begin{align}
        \bigl\lVert S_{\Gamma_n^{\mathrm{ES}}}
        -\mathcal K_\infty^{-1}\bigr\rVert
         & \leq
        \frac{a}{1-a}\,\bigl\lVert\mathcal K_\infty^{-1}\bigr\rVert.
    \end{align}
\end{theorem}

\begin{proof}
    \leanok
    \uses{thm:primitive_ground_space_gram_eventual_inverse_bound, thm:inverse_gram_identities}
    Eventual invertibility of
    $(\Gamma_n^{\mathrm{ES}})^\dagger\Gamma_n^{\mathrm{ES}}$ implies eventual
    injectivity of $\Gamma_n^{\mathrm{ES}}$.  The inverse-Gram identities then
    identify $S_{\Gamma_n^{\mathrm{ES}}}$ with $\mathcal K_n^{-1}$, and the
    preceding estimate applies.
\end{proof}

\begin{definition}[Injective range-projector candidate]
    \label{def:injective_range_projector}
    \lean{ContinuousLinearMap.injectiveRangeProjector}
    \leanok
    \uses{def:inverse_gram}
    For injective $T$, define the operator $Q_T:=T S_T T^\dagger$ on its
    codomain.
\end{definition}

\begin{theorem}[Injective range-projector formula]
    \label{thm:injective_range_projector_inverse_gram}
    \lean{ContinuousLinearMap.injectiveRangeProjector_eq_starProjection}
    \leanok
    \uses{def:injective_range_projector}
    If $T$ is an injective map from a finite-dimensional Hilbert space into a
    complete Hilbert space, then
    \begin{align}
        P_{\range T} & =T(T^\dagger T)^{-1}T^\dagger.
    \end{align}
\end{theorem}

\begin{proof}
    \leanok
    \uses{def:inverse_gram, thm:inverse_gram_identities}
    The vector $Q_Tx$ lies in $\range T$.  Moreover,
    $T^\dagger Q_Tx=(T^\dagger T)S_TT^\dagger x=T^\dagger x$, so for every
    $y$ the adjoint identity gives
    $\langle x-Q_Tx,Ty\rangle=0$.  Thus $x-Q_Tx$ is orthogonal to
    $\range T$, and the range--orthogonal decomposition characterizing the
    orthogonal projector gives $Q_T=P_{\range T}$.
\end{proof}

% Maintainer note: this is the dependency-closed projector slice for issue #5752.
% Its source boundary is the inverse-Gram/range-projector identity; it does not
% include or assert the FNW overlapping-window estimate.
\begin{theorem}[Physical ground-space range projector]
    \label{thm:ground_space_map_es_injective_range_projector}
    \lean{MPSTensor.injectiveRangeProjector_groundSpaceMapES_eq_starProjection}
    \leanok
    \uses{def:ground_space_map_es, def:ground_space_es, def:injective_range_projector, thm:ground_space_map_es_injective_block}
    If $A$ is $L$-block injective, then the injective range projector of
    $\Gamma_L^{\mathrm{ES}}$ is the physical orthogonal projector onto the
    local ground space:
    \begin{align}
        Q_{\Gamma_L^{\mathrm{ES}}}
         & =P_{G_L^{\mathrm{ES}}(A)}.
    \end{align}
    In particular, this is the physical range projector on the local Hilbert
    space, not the virtual transfer fixed-point projector.
\end{theorem}

\begin{proof}
    \leanok
    \uses{thm:injective_range_projector_inverse_gram, thm:range_ground_space_map_es}
    The general injective range-projector theorem identifies
    $Q_{\Gamma_L^{\mathrm{ES}}}$ with the orthogonal projector onto
    $\range(\Gamma_L^{\mathrm{ES}})$.  The exact-range theorem identifies
    this range with $G_L^{\mathrm{ES}}(A)$.
\end{proof}

\begin{theorem}[Physical ground-space inverse-Gram projector formula]
    \label{thm:ground_space_es_projection_inverse_gram}
    \lean{MPSTensor.groundSpaceES_starProjection_eq_groundSpaceMapES_comp_inverseGram_comp_adjoint}
    \leanok
    \uses{def:ground_space_es, def:ground_space_map_es, def:inverse_gram, thm:ground_space_map_es_injective_block_of_le}
    Suppose that $L_0>0$, that $A$ is $L_0$-block injective, and that
    $L\ge L_0$.  Then the physical orthogonal projector onto the local ground
    space is
    \begin{align}
        P_{G_L^{\mathrm{ES}}(A)}
         & =\Gamma_L^{\mathrm{ES}}
        \left((\Gamma_L^{\mathrm{ES}})^\dagger
        \Gamma_L^{\mathrm{ES}}\right)^{-1}
        (\Gamma_L^{\mathrm{ES}})^\dagger.
    \end{align}
    Thus the displayed operator is the physical range projector on
    $\ell^2(\Fin{d}^{L})$.
\end{theorem}

\begin{proof}
    \leanok
    \uses{thm:ground_space_map_es_injective_range_projector, thm:wordspan_blk_inj_succ}
    Positive-length propagation makes $A$ block injective at length $L$.
    Apply the physical range-projector identity at that length and unfold its
    inverse-Gram expression.
\end{proof}

\begin{theorem}[Neumann inverse estimate]
    \label{thm:neumann_inverse_norm_bound}
    \lean{NormedRing.norm_inverse_one_sub_le}
    \leanok
    If $\lVert t\rVert\le a<1$ in a normed ring with convergent geometric
    series and $\lVert 1\rVert=1$, then
    \begin{align}
        \lVert(1-t)^{-1}\rVert & \le(1-a)^{-1}.
    \end{align}
\end{theorem}

\begin{proof}
    \leanok
    Since $\lVert t\rVert\le a<1$, the Neumann series converges and
    $(1-t)^{-1}=\sum_{n\ge0}t^n$.  The triangle inequality and
    submultiplicativity give
    \begin{align}
        \left\lVert\sum_{n\ge0}t^n\right\rVert
         & \le \sum_{n\ge0}\lVert t\rVert^n
        =(1-\lVert t\rVert)^{-1}
        \le(1-a)^{-1},
    \end{align}
    where the last step uses monotonicity of the reciprocal on the positive
    reals.
\end{proof}

\begin{theorem}[Quantitative inverse-Gram estimate]
    \label{thm:inverse_gram_norm_bound}
    \lean{ContinuousLinearMap.norm_inverseGram_le_of_norm_one_sub_adjoint_comp_self_le}
    \leanok
    \uses{def:inverse_gram}
    Let $\mathbb K$ be a real or complex scalar field, let $E$ be a
    finite-dimensional complete $\mathbb K$-Hilbert space, and let $F$ be a
    complete $\mathbb K$-Hilbert space.  Let $T\colon E\to F$ be an injective
    continuous linear map.  If $a\in\mathbb R$ satisfies $a<1$ and
    \begin{align}
        \lVert I_E-T^\dagger T\rVert & \le a,
    \end{align}
    then
    \begin{align}
        \lVert S_T\rVert & \le(1-a)^{-1}.
    \end{align}
\end{theorem}

\begin{proof}
    \leanok
    \uses{thm:inverse_gram_identities, thm:neumann_inverse_norm_bound}
    If $E$ is trivial, then $S_T=0$ and the bound follows from $a<1$.
    Otherwise, identify $S_T$ with the ring inverse of $T^\dagger T$ and apply
    Theorem~\ref{thm:neumann_inverse_norm_bound} to
    $t=I_E-T^\dagger T$.
\end{proof}

\begin{definition}[Canonical parent interaction]\label{def:parent_interaction}
    \lean{MPSTensor.parentInteraction}
    \leanok
    \uses{def:ground_space_parent}
    The parent interaction at range $L$ is the canonical orthogonal projector
    specified in (\ref{eq:parent_interaction}).
\end{definition}

\begin{lemma}[Parent interaction is idempotent]
    \label{lem:parent_interaction_idempotent}
    \lean{MPSTensor.parentInteraction_idempotent}
    \leanok
    \uses{def:parent_interaction}
    The local parent interaction is a projector:
    $h_L(A)^2=h_L(A)$.
\end{lemma}

\begin{proof}
    \leanok
    This is the idempotency of the orthogonal projector onto $G_L(A)^\perp$,
    transported from the Hilbert-space realization of the local coefficient
    space.
\end{proof}

\begin{lemma}[Parent interaction annihilates ground-space elements]%
    \label{lem:parent_interaction_apply_mem_ground_space}
    \lean{MPSTensor.parentInteraction_apply_mem_groundSpace}
    \leanok
    \uses{def:parent_interaction, def:ground_space_parent}
    For any $v\in G_L(A)$, one has $h_L(A)v=0$.
\end{lemma}

\begin{proof}
    \leanok
    Since $h_L(A)$ is the orthogonal projector onto $G_L(A)^\perp$, it
    annihilates every element of $G_L(A)$.
\end{proof}

\begin{definition}[Restriction to a cyclic interval]\label{def:extract_window}
    \lean{MPSTensor.extractWindow}
    \leanok
    \uses{def:nsite_space_parent}
    For $\sigma\in\Fin{d}^{N}$ and a starting site $i$, write
    $\sigma|_{[i,i+L-1]}\in\Fin{d}^{L}$ for the word defined by
    $(\sigma|_{[i,i+L-1]})(r)=\sigma(i+r\bmod N)$ for $0\le r<L$.
\end{definition}

\begin{definition}[Prescribing a cyclic interval]\label{def:replace_window}
    \lean{MPSTensor.replaceWindow}
    \leanok
    \uses{def:extract_window}
    If $L\le N$ and $\tau\in\Fin{d}^{L}$, write
    $\sigma^{[i,i+L-1]\leftarrow\tau}$ for the configuration defined by
    \begin{align}
        \sigma^{[i,i+L-1]\leftarrow\tau}(k)
         & =
        \begin{cases}
            \tau(r),
             & k\equiv i+r\pmod N
            \text{ for some }0\le r<L, \\
            \sigma(k),
             & \text{otherwise}.
        \end{cases}
        \notag
    \end{align}
    The index $r$ in the first case is unique because $L\le N$.
\end{definition}

\begin{lemma}[Extracting a replaced window]\label{lem:extract_window_replace_window}
    \lean{MPSTensor.extractWindow_replaceWindow}
    \leanok
    \uses{def:extract_window, def:replace_window}
    If $L\le N$, then restricting to the same cyclic interval after prescribing
    its values recovers the inserted word:
    $\left(\sigma^{[i,i+L-1]\leftarrow\tau}\right)|_{[i,i+L-1]}=\tau$.
\end{lemma}

\begin{lemma}[Replacing an extracted window]\label{lem:replace_window_extract_window}
    \lean{MPSTensor.replaceWindow_extractWindow}
    \leanok
    \uses{def:extract_window, def:replace_window}
    If $L\le N$, then prescribing on a cyclic interval the values already
    carried by $\sigma$ leaves $\sigma$ unchanged:
    $\sigma^{[i,i+L-1]\leftarrow\sigma|_{[i,i+L-1]}}=\sigma$.
\end{lemma}

\begin{proof}
    \leanok
    Case-split on whether the offset $(k-i+N)\bmod N$ lies in $[0,L)$:
    inside the cyclic interval the prescribed value is the original value of
    $\sigma$; outside the interval both sides already agree with $\sigma$.
\end{proof}

\begin{lemma}[Replacing the same window twice]\label{lem:replace_window_replace_window_same}
    \lean{MPSTensor.replaceWindow_replaceWindow_same}
    \leanok
    \uses{def:replace_window}
    If $L\le N$, then two prescriptions on the same cyclic interval reduce to
    the second prescription:
    \begin{align}
        \left(\sigma^{[i,i+L-1]\leftarrow\tau}\right)^{[i,i+L-1]\leftarrow\upsilon}
         & =
        \sigma^{[i,i+L-1]\leftarrow\upsilon}.
        \notag
    \end{align}
\end{lemma}

\begin{proof}
    \leanok
    On the cyclic interval both sides take the value prescribed by $\upsilon$.
    Away from that interval neither second prescription changes the value.
\end{proof}

For a periodic chain of length $N$, the parent Hamiltonian is the sum of
translated copies of the local interaction.

\begin{definition}[Translated local term]\label{def:local_term_parent}
    \lean{MPSTensor.localTerm}
    \leanok
    \uses{def:parent_interaction, def:replace_window, def:extract_window}
    For each site index $i\in\Fin{N}$, the translated local term is determined
    by
    \begin{align}
        (h_i\psi)(\sigma)
         & =
        \begin{cases}
            \left[h_L(A)
                \left(\tau\mapsto
                \psi\left(\sigma^{[i,i+L-1]\leftarrow\tau}\right)\right)\right]
            \left(\sigma|_{[i,i+L-1]}\right),
             & L\le N, \\
            0,
             & L>N.
        \end{cases}
        \label{eq:parent_local_term}
    \end{align}
\end{definition}

\begin{lemma}[Pointwise form of a translated local term]\label{lem:local_term_apply_of_le}
    \lean{MPSTensor.localTerm_apply_of_le}
    \leanok
    \uses{def:local_term_parent}
    If $L\le N$, then the translated term is obtained by applying the local
    parent interaction on the chosen cyclic window:
    \begin{align}
        (h_i(A,L)\psi)(\sigma)
         & =
        \left[h_L(A)
            \left(\tau\mapsto
            \psi\left(\sigma^{[i,i+L-1]\leftarrow\tau}\right)\right)\right]
        \left(\sigma|_{[i,i+L-1]}\right).
        \notag
    \end{align}
\end{lemma}

\begin{proof}
    \leanok
    This is the first case in (\ref{eq:parent_local_term}).
\end{proof}

\begin{lemma}[Translated local terms are idempotent]\label{lem:local_term_idempotent}
    \lean{MPSTensor.localTerm_idempotent}
    \leanok
    \uses{def:local_term_parent, lem:parent_interaction_idempotent, lem:local_term_apply_of_le, lem:extract_window_replace_window, lem:replace_window_replace_window_same}
    For every site $i$, one has $h_i(A,L)^2=h_i(A,L)$.
\end{lemma}

\begin{proof}
    \leanok
    If $L>N$, the translated term is zero. Suppose $L\le N$ and put
    $f_{\sigma,i}(\tau)
        :=\psi(\sigma^{[i,i+L-1]\leftarrow\tau})$.
    The pointwise formula for translated terms, together with
    Lemmas~\ref{lem:extract_window_replace_window} and
    \ref{lem:replace_window_replace_window_same}, gives
    \begin{align}
        (h_i(A,L)^2\psi)(\sigma)
         & =
        (h_L(A)^2f_{\sigma,i})
        \left(\sigma|_{[i,i+L-1]}\right) =
        (h_L(A)f_{\sigma,i})
        \left(\sigma|_{[i,i+L-1]}\right) =
        (h_i(A,L)\psi)(\sigma).
        \notag
    \end{align}
    Here the second equality follows from
    Lemma~\ref{lem:parent_interaction_idempotent}.
\end{proof}

\begin{definition}[Parent Hamiltonian]\label{def:parent_hamiltonian}
    \lean{MPSTensor.parentHamiltonian}
    \leanok
    \uses{def:local_term_parent}
    The chain Hamiltonian is
    \begin{align}
        H_N(A,L)
         & :=\sum_{i=0}^{N-1}h_i.
        \label{eq:parent_hamiltonian}
    \end{align}
\end{definition}

\begin{definition}[Frustration-free condition]\label{def:is_frustration_free}
    \lean{MPSTensor.IsFrustrationFree}
    \leanok
    \uses{def:local_term_parent}
    A vector $\psi$ is a frustration-free ground state of the parent Hamiltonian
    when each local interaction annihilates it, that is, when
    $h_i\psi=0$ for every $i\in\Fin{N}$.
\end{definition}

\begin{lemma}[Periodic MPS vector window membership]\label{lem:mpv_window_mem_ground_space}
    \lean{MPSTensor.mpv_window_mem_groundSpace}
    \leanok
    \uses{def:ground_space_parent, def:mpv}
    For any window of $L$ consecutive sites starting at position~$i$ in an
    $N$-site periodic chain with $L\le N$, the restriction of the periodic MPS
    vector $V^{(N)}(A)$ to that window lies in $G_L(A)$.
\end{lemma}

\begin{proof}
    \leanok
    The boundary matrix is the product of the $A$-matrices on the complement
    sites. Trace cyclicity rotates the full $N$-site product so that the window
    indices come first, matching the definition of the ground-space map.
\end{proof}

\begin{lemma}[Local term annihilates the periodic MPS vector]\label{lem:local_term_annihilates_mpv}
    \lean{MPSTensor.localTerm_annihilates_mpv}
    \leanok
    \uses{def:local_term_parent, def:mpv, lem:mpv_window_mem_ground_space, lem:parent_interaction_apply_mem_ground_space}
    For each site $i$, one has $h_iV^{(N)}(A)=0$.
\end{lemma}

\begin{proof}
    \leanok
    This combines periodic-MPS-vector window membership with the fact that the
    parent interaction annihilates ground-space elements.
\end{proof}

\begin{lemma}[Parent Hamiltonian annihilates the periodic MPS vector]\label{lem:parent_hamiltonian_annihilates}
    \lean{MPSTensor.parentHamiltonian_annihilates}
    \leanok
    \uses{def:parent_hamiltonian, def:mpv, lem:local_term_annihilates_mpv}
    For every $N$ with $L\le N$, the periodic MPS vector satisfies
    $H_N(A,L)V^{(N)}(A)=0$.
\end{lemma}

\begin{proof}
    \leanok
    This is a sum of zeros, since each local term annihilates the periodic MPS
    vector.
\end{proof}

\begin{lemma}[Frustration-freeness of the periodic MPS vector]\label{lem:parent_hamiltonian_ff}
    \lean{MPSTensor.parentHamiltonian_frustrationFree}
    \leanok
    \uses{def:is_frustration_free, def:mpv, lem:local_term_annihilates_mpv}
    The periodic MPS vector is a frustration-free ground state: for every site
    $i$, one has $h_iV^{(N)}(A)=0$.
\end{lemma}

\begin{proof}
    \leanok
    The claim follows from Lemma~\ref{lem:local_term_annihilates_mpv}.
\end{proof}

\section{Intersection property and unique ground state}\label{sec:intersection_unique}

This section develops the intersection property of local ground spaces for
injective MPS and the resulting uniqueness of the ground state on the periodic
chain.

\subsection{Restriction maps}

Given a state $\psi$ on $L+1$ sites, we define two restriction maps.

\begin{remark}[Word and restriction identities]
    \lean{List.ofFn_snoc}
    \lean{MPSTensor.evalWord_ofFn_snoc}
    \lean{MPSTensor.evalWord_ofFn_cons}
    \lean{MPSTensor.restrictLastₗ}
    \lean{MPSTensor.restrictFirstₗ}
    \lean{MPSTensor.restrictLast_apply}
    \lean{MPSTensor.restrictFirst_apply}
    \leanok
    Appending a last letter to a word corresponds to right multiplication by the
    corresponding tensor letter, while prepending a first letter corresponds to
    left multiplication. The restriction maps are the associated linear maps on
    the state spaces under the computational-basis identification, together with
    their pointwise formulas.
\end{remark}

\begin{definition}[Left restriction]\label{def:restrict_last}
    \lean{MPSTensor.restrictLast}
    \leanok
    Fixing the last physical index $j$ gives the state
    $\sigma\mapsto\psi(\sigma_1,\ldots,\sigma_L,j)$.
\end{definition}

\begin{definition}[Right restriction]\label{def:restrict_first}
    \lean{MPSTensor.restrictFirst}
    \leanok
    Fixing the first physical index $i$ gives the state
    $\sigma\mapsto\psi(i,\sigma_1,\ldots,\sigma_L)$.
\end{definition}

\begin{definition}[Left ground membership]\label{def:in_left_ground}
    \lean{MPSTensor.InLeftGround}
    \leanok
    \uses{def:restrict_last, def:ground_space_parent}
    A state $\psi$ on $L+1$ sites satisfies the left ground condition if, for
    every $j$, the state obtained by fixing the last index to $j$ lies in
    $G_L(A)$.
\end{definition}

\begin{definition}[Right ground membership]\label{def:in_right_ground}
    \lean{MPSTensor.InRightGround}
    \leanok
    \uses{def:restrict_first, def:ground_space_parent}
    A state $\psi$ on $L+1$ sites satisfies the right ground condition if, for
    every $i$, the state obtained by fixing the first index to $i$ lies in
    $G_L(A)$.
\end{definition}

\subsection{Forward direction: ground space restricts}

\begin{lemma}[Left restriction]\label{lem:ground_space_in_left_ground}
    \lean{MPSTensor.groundSpace_inLeftGround}
    \leanok
    \uses{def:ground_space_parent, def:in_left_ground}
    If $\psi\in G_{L+1}(A)$, i.e.\ if
    $\psi(\sigma)=\tr(A^\sigma X)$ for some~$X$, then, for each $j$, the state
    obtained by fixing the last index to $j$ lies in $G_L(A)$ with witness
    $A^jX$.
\end{lemma}

\begin{proof}
    \leanok
    The restriction simply absorbs the last tensor into the boundary matrix:
    \begin{center}
        \tenkzkernel
        % Formula/source: Papers/quant-ph_0608197/MPSarchive.tex,
        % lines 889--898, with boundary argument A^j X in Gamma_L.
        % Ink-to-index map: sigma_1,...,sigma_L remain free; j is absorbed into
        % the boundary capsule and carries no free leg in this shortened state.
        % Contraction set: every virtual bond of the periodic A-word--A^jX loop.
        % Type verdict: a rank-L physical tensor in G_L(A).
        % Boundary signature: semantic virtual-west=0 | virtual-east=0 |
        % physical-up=L | physical-down=0; the finite schematic emits 3 up legs.
        \begin{tenkz}[rows={wire}, cols=5, boundary=periodic]
            \tn[skin=ring]{A^j X} &
            \tn[label pos=135, ports={90:physical:$\sigma_1$}]{A} &
            \tn[label pos=135, ports={90:physical}]{A} & \tn[skin=dots]{} &
            \tn[label pos=135, ports={90:physical:$\sigma_L$}]{A}
            \tnmark[form=bracket]{(1,2) .. (1,5)}{$L$}
        \end{tenkz}
    \end{center}
    Direct computation gives
    \begin{align}
        \psi(\sigma_1,\ldots,\sigma_L,j)
         & =
        \tr\!\left(A^{\sigma_1}\cdots A^{\sigma_L}A^jX\right)
        =
        \Gamma_L(A^jX)(\sigma).
        \notag
    \end{align}
\end{proof}

\begin{lemma}[Right restriction]\label{lem:ground_space_in_right_ground}
    \lean{MPSTensor.groundSpace_inRightGround}
    \leanok
    \uses{def:ground_space_parent, def:in_right_ground}
    If $\psi\in G_{L+1}(A)$ with $\psi(\sigma)=\tr(A^\sigma X)$, then, for each
    $i$, the state obtained by fixing the first index to $i$ lies in $G_L(A)$
    with witness $XA^i$.
\end{lemma}

\begin{proof}
    \leanok
    On the right restriction, one first rotates the trace and then reads the
    resulting boundary matrix on the shortened window:
    \begin{center}
        \tenkzkernel
        % Formula/source: Papers/quant-ph_0608197/MPSarchive.tex,
        % lines 889--898, with the cyclically rotated boundary argument X A^i.
        % Ink-to-index map: sigma_1,...,sigma_L remain free; i has been absorbed
        % into the boundary capsule and is not a free leg of the shortened state.
        % Contraction set: every virtual bond of the periodic A-word--XA^i loop.
        % Type verdict: a rank-L physical tensor in G_L(A).
        % Boundary signature: semantic virtual-west=0 | virtual-east=0 |
        % physical-up=L | physical-down=0; the finite schematic emits 3 up legs.
        \begin{tenkz}[rows={wire}, cols=5, boundary=periodic]
            \tn[skin=ring]{X A^i} &
            \tn[label pos=135, ports={90:physical:$\sigma_1$}]{A} &
            \tn[label pos=135, ports={90:physical}]{A} & \tn[skin=dots]{} &
            \tn[label pos=135, ports={90:physical:$\sigma_L$}]{A}
            \tnmark[form=bracket]{(1,2) .. (1,5)}{$L$}
        \end{tenkz}
    \end{center}
    By trace cyclicity,
    \begin{align}
        \psi(i,\sigma_1,\ldots,\sigma_L)
         & =
        \tr(A^iA^\sigma X)
        =
        \tr(A^\sigma XA^i)
        =
        \Gamma_L(XA^i)(\sigma).
        \notag
    \end{align}
\end{proof}
